\documentclass[11pt,a4paper,openany]{book}

% =============================================================================
% Professional textbook foundation
% =============================================================================

\usepackage[T1]{fontenc}
\usepackage[utf8]{inputenc}
\usepackage{lmodern}
\usepackage{microtype}
\usepackage[a4paper,margin=27mm,headheight=15pt]{geometry}
\usepackage{setspace}
\setstretch{1.05}
\setlength{\parindent}{1.15em}
\setlength{\parskip}{0pt}
\raggedbottom

\usepackage{amsmath,amssymb,mathtools}
\usepackage{booktabs,tabularx,array}

\usepackage{tikz}
\usepackage{pgfplots}
\usepackage{caption}
\usepackage{subcaption}
\usepackage{float}
\usepackage{wrapfig}
\usepackage{adjustbox}
\usetikzlibrary{
  arrows.meta,
  angles,
  quotes,
  calc,
  positioning,
  intersections,
  decorations.pathmorphing,
  decorations.markings,
  backgrounds,
  fit,
  shapes.geometric,
  patterns,
  3d
}
\pgfplotsset{compat=1.18}

\newcolumntype{Y}{>{\raggedright\arraybackslash}X}

\usepackage{xcolor}
\definecolor{chapterblue}{RGB}{35,75,120}
\definecolor{softblue}{RGB}{240,246,252}
\definecolor{softgreen}{RGB}{241,249,244}
\definecolor{softgold}{RGB}{252,248,236}
\definecolor{softred}{RGB}{252,242,242}
\definecolor{softgray}{RGB}{246,247,249}
\definecolor{rulegray}{RGB}{185,190,198}

\usepackage[most]{tcolorbox}
\tcbset{
  enhanced,
  breakable,
  boxrule=0.6pt,
  arc=1.4mm,
  left=2mm,right=2mm,top=1.5mm,bottom=1.5mm,
  before skip=7pt,after skip=7pt
}

\newtcolorbox{keyidea}[1][Key Idea]{
  title=#1,colback=softblue,colframe=chapterblue,fonttitle=\bfseries
}
\newtcolorbox{definitionbox}[1][Definition]{
  title=#1,colback=softgreen,colframe=green!45!black,fonttitle=\bfseries
}
\newtcolorbox{physicalbox}[1][Physical Interpretation]{
  title=#1,colback=softgold,colframe=orange!55!black,fonttitle=\bfseries
}
\newtcolorbox{mistakebox}[1][Common Mistake]{
  title=#1,colback=softred,colframe=red!55!black,fonttitle=\bfseries
}
\newtcolorbox{historybox}[1][Historical Note]{
  title=#1,colback=softgray,colframe=rulegray,fonttitle=\bfseries
}
\newtcolorbox{examplebox}[1][Worked Example]{
  title=#1,colback=white,colframe=chapterblue,fonttitle=\bfseries
}
\newtcolorbox{solutionbox}[1][Solution]{
  title=#1,colback=softgray,colframe=rulegray,fonttitle=\bfseries
}
\newtcolorbox{roadmapbox}[1][Chapter Roadmap]{
  title=#1,colback=softgray,colframe=chapterblue,fonttitle=\bfseries
}

\usepackage{enumitem}
\setlist[itemize]{leftmargin=1.8em,topsep=4pt,itemsep=2pt,parsep=0pt,partopsep=0pt}
\setlist[enumerate]{leftmargin=2em,topsep=4pt,itemsep=3pt,parsep=0pt}

\usepackage{titlesec}
\titleformat{\chapter}[display]
  {\normalfont\bfseries\color{chapterblue}}
  {\filleft\Large\chaptername\ \thechapter}
  {0.7ex}
  {\titlerule\vspace{1.1ex}\Huge}
  [\vspace{0.7ex}\titlerule]
\titlespacing*{\chapter}{0pt}{-16pt}{20pt}

\titleformat{\section}
  {\Large\bfseries\color{chapterblue}}
  {\thesection}{0.7em}{}
\titlespacing*{\section}{0pt}{2.1ex plus .5ex minus .2ex}{0.7ex}

\titleformat{\subsection}
  {\large\bfseries}
  {\thesubsection}{0.7em}{}
\titlespacing*{\subsection}{0pt}{1.7ex plus .4ex minus .2ex}{0.5ex}

\usepackage{fancyhdr}
\pagestyle{fancy}
\fancyhf{}
\fancyhead[LE]{\small\itshape Theory of Quantum Mechanics}
\fancyhead[RO]{\small\itshape\nouppercase{\rightmark}}
\fancyfoot[C]{\thepage}
\renewcommand{\headrulewidth}{0.4pt}
\renewcommand{\footrulewidth}{0pt}
\renewcommand{\chaptermark}[1]{\markboth{\thechapter.\ #1}{}}
\renewcommand{\sectionmark}[1]{\markright{\thesection.\ #1}}

\usepackage{tocloft}
\setcounter{tocdepth}{2}
\setcounter{secnumdepth}{3}
\renewcommand{\contentsname}{Table of Contents}
\renewcommand{\cftchapfont}{\bfseries}
\renewcommand{\cftchappagefont}{\bfseries}
\setlength{\cftbeforechapskip}{7pt}
\setlength{\cftbeforesecskip}{1pt}
\renewcommand{\cftdotsep}{1.5}

\usepackage[hidelinks]{hyperref}

\newcommand{\chapterquote}[2]{%
  \begin{center}
    \begin{minipage}{0.86\textwidth}
      \centering\itshape\small #1
      \if\relax\detokenize{#2}\relax
      \else
        \par\vspace{0.45em}\raggedleft\normalfont\footnotesize--- #2
      \fi
    \end{minipage}
  \end{center}
  \vspace{0.7em}
}

\newcommand{\learningobjectives}{%
  \section*{Learning Objectives}
  \addcontentsline{toc}{section}{Learning Objectives}
  \noindent After completing this chapter, the reader should be able to:
}

\newcommand{\chaptersummary}{%
  \section*{Chapter Summary}
  \addcontentsline{toc}{section}{Chapter Summary}
}

\newcommand{\nextchapter}{%
  \section*{Looking Ahead}
  \addcontentsline{toc}{section}{Looking Ahead}
}


%%%%%%%%%%%%%%%%%%%%%%%%%%%%%%%%%%%%%%%%%%%%%%%%%%%%%%%%%%%%%%%%%%%%%%%%%%%%%%%%
% PR-BOOK-02 : Professional textbook environments
%%%%%%%%%%%%%%%%%%%%%%%%%%%%%%%%%%%%%%%%%%%%%%%%%%%%%%%%%%%%%%%%%%%%%%%%%%%%%%%%

\newcounter{definition}[chapter]
\renewcommand{\thedefinition}{\thechapter.\arabic{definition}}

\newcounter{example}[chapter]
\renewcommand{\theexample}{\thechapter.\arabic{example}}

\newcounter{remark}[chapter]
\renewcommand{\theremark}{\thechapter.\arabic{remark}}

\newenvironment{definition}{
\refstepcounter{definition}
\begin{definitionbox}[Definition \thedefinition]
}{
\end{definitionbox}
}

\newenvironment{workedexample}{
\refstepcounter{example}
\begin{examplebox}[Example \theexample]
}{
\end{examplebox}
}

\newenvironment{solution}{
\begin{solutionbox}[Solution]
}{
\end{solutionbox}
}

\newenvironment{physical}{
\begin{physicalbox}[Physical Interpretation]
}{
\end{physicalbox}
}

\newenvironment{warningbox}{
\begin{mistakebox}[Common Mistake]
}{
\end{mistakebox}
}

\newenvironment{historicalnote}{
\begin{historybox}[Historical Note]
}{
\end{historybox}
}

\newcommand{\ImportantEquation}[1]{
\begin{keyidea}[Key Equation]
\[
#1
\]
\end{keyidea}
}



%%%%%%%%%%%%%%%%%%%%%%%%%%%%%%%%%%%%%%%%%%%%%%%%%%%%%%%%%%%%%%%%%%%%%%%%%%%%%%%%
% PR-BOOK-03 : Complete educational framework
%%%%%%%%%%%%%%%%%%%%%%%%%%%%%%%%%%%%%%%%%%%%%%%%%%%%%%%%%%%%%%%%%%%%%%%%%%%%%%%%

% Additional counters ---------------------------------------------------------
\newcounter{keyequation}[chapter]
\renewcommand{\thekeyequation}{\thechapter.\arabic{keyequation}}

\newcounter{exercise}[chapter]
\renewcommand{\theexercise}{\thechapter.\arabic{exercise}}

\newcounter{historicalnote}[chapter]
\renewcommand{\thehistoricalnote}{\thechapter.\arabic{historicalnote}}

% Additional colours ----------------------------------------------------------
\definecolor{summaryblue}{RGB}{232,241,250}
\definecolor{notationgray}{RGB}{248,249,251}
\definecolor{exercisegreen}{RGB}{239,248,241}
\definecolor{outcomegold}{RGB}{253,248,229}
\definecolor{lookaheadviolet}{RGB}{247,242,251}

% Educational boxes -----------------------------------------------------------
\newtcolorbox{keyequationbox}[1]{
  title={Key Equation \thekeyequation: #1},
  colback=softblue,
  colframe=chapterblue,
  fonttitle=\bfseries,
  borderline west={2.2pt}{0pt}{chapterblue}
}

\newtcolorbox{summarybox}[1][Chapter Summary]{
  title=#1,
  colback=summaryblue,
  colframe=chapterblue,
  fonttitle=\bfseries\large
}

\newtcolorbox{formulasheetbox}[1][Quick Formula Sheet]{
  title=#1,
  colback=white,
  colframe=chapterblue,
  fonttitle=\bfseries\large,
  boxrule=0.8pt
}

\newtcolorbox{notationbox}[1][Notation and Units]{
  title=#1,
  colback=notationgray,
  colframe=rulegray,
  fonttitle=\bfseries\large
}

\newtcolorbox{outcomesbox}[1][Learning Outcomes Check]{
  title=#1,
  colback=outcomegold,
  colframe=orange!60!black,
  fonttitle=\bfseries\large
}

\newtcolorbox{lookingaheadbox}[1][Looking Ahead]{
  title=#1,
  colback=lookaheadviolet,
  colframe=violet!55!black,
  fonttitle=\bfseries\large
}

\newtcolorbox{exercisebox}[2][]{
  title={Exercise \theexercise\if\relax\detokenize{#2}\relax\else: #2\fi},
  colback=exercisegreen,
  colframe=green!45!black,
  fonttitle=\bfseries,
  #1
}

\newtcolorbox{dashboardbox}[1][Chapter Dashboard]{
  title=#1,
  colback=softgray,
  colframe=rulegray,
  fonttitle=\bfseries\large
}

% Semantic environments -------------------------------------------------------
\newenvironment{keyequation}[1]{%
  \refstepcounter{keyequation}%
  \begin{keyequationbox}{#1}%
}{%
  \end{keyequationbox}%
}

\newenvironment{chapterexercise}[1][]{%
  \refstepcounter{exercise}%
  \begin{exercisebox}{#1}%
}{%
  \end{exercisebox}%
}

\newenvironment{chapterhistorical}[1][]{%
  \refstepcounter{historicalnote}%
  \begin{historybox}[Historical Note \thehistoricalnote%
  \if\relax\detokenize{#1}\relax\else: #1\fi]%
}{%
  \end{historybox}%
}

% Structured worked-example helpers ------------------------------------------
\newcommand{\problempart}{\par\medskip\noindent\textbf{Problem.}\ }
\newcommand{\givenpart}{\par\medskip\noindent\textbf{Given.}\ }
\newcommand{\findpart}{\par\medskip\noindent\textbf{Find.}\ }
\newcommand{\solutionpart}{\par\medskip\noindent\textbf{Solution.}\ }
\newcommand{\verificationpart}{\par\medskip\noindent\textbf{Verification.}\ }
\newcommand{\meaningpart}{\par\medskip\noindent\textbf{Physical meaning.}\ }

% Difficulty tags -------------------------------------------------------------
\newcommand{\difficulty}[1]{%
  \hfill\textbf{Difficulty: }%
  \ifcase#1\or$\star\circ\circ\circ\circ$%
  \or$\star\star\circ\circ\circ$%
  \or$\star\star\star\circ\circ$%
  \or$\star\star\star\star\circ$%
  \or$\star\star\star\star\star$\fi
}

% Consistent chapter-ending headings ------------------------------------------
\newcommand{\chapterformulaheading}{%
  \section*{Quick Formula Sheet}
  \addcontentsline{toc}{section}{Quick Formula Sheet}
}
\newcommand{\chapternotationheading}{%
  \section*{Notation and Units}
  \addcontentsline{toc}{section}{Notation and Units}
}
\newcommand{\chapteroutcomesheading}{%
  \section*{Learning Outcomes Check}
  \addcontentsline{toc}{section}{Learning Outcomes Check}
}
\newcommand{\chapterexercisesheading}{%
  \section*{Exercises}
  \addcontentsline{toc}{section}{Exercises}
}
\newcommand{\chapterlookaheadheading}{%
  \section*{Looking Ahead}
  \addcontentsline{toc}{section}{Looking Ahead}
}

% Compact reusable tables -----------------------------------------------------
\newenvironment{notationtable}{%
  \begin{tabular}{>{\centering\arraybackslash}p{0.16\textwidth}
                  >{\raggedright\arraybackslash}p{0.48\textwidth}
                  >{\raggedright\arraybackslash}p{0.23\textwidth}}
  \toprule
  \textbf{Symbol} & \textbf{Meaning} & \textbf{SI unit}\\
  \midrule
}{%
  \bottomrule
  \end{tabular}
}

\newenvironment{formulatable}{%
  \begin{tabular}{>{\raggedright\arraybackslash}p{0.29\textwidth}
                  >{\raggedright\arraybackslash}p{0.61\textwidth}}
  \toprule
  \textbf{Quantity or operation} & \textbf{Formula}\\
  \midrule
}{%
  \bottomrule
  \end{tabular}
}



%%%%%%%%%%%%%%%%%%%%%%%%%%%%%%%%%%%%%%%%%%%%%%%%%%%%%%%%%%%%%%%%%%%%%%%%%%%%%%%%
% PR-BOOK-04 : Visual identity and TikZ illustration system
%%%%%%%%%%%%%%%%%%%%%%%%%%%%%%%%%%%%%%%%%%%%%%%%%%%%%%%%%%%%%%%%%%%%%%%%%%%%%%%%

% Visual palette ----------------------------------------------------------------
\definecolor{visualblue}{RGB}{35,75,120}
\definecolor{visualcyan}{RGB}{54,135,170}
\definecolor{visualgold}{RGB}{193,139,46}
\definecolor{visualred}{RGB}{164,63,63}
\definecolor{visualgreen}{RGB}{58,126,86}
\definecolor{visualink}{RGB}{38,45,54}
\definecolor{visualgrid}{RGB}{210,218,226}

% Global TikZ styles ------------------------------------------------------------
\tikzset{
  physics axis/.style={
    -{Latex[length=3mm]},
    line width=0.9pt,
    draw=visualink
  },
  physics vector/.style={
    -{Latex[length=3.2mm]},
    line width=1.5pt,
    draw=visualblue
  },
  physics vector secondary/.style={
    -{Latex[length=3mm]},
    line width=1.25pt,
    draw=visualcyan
  },
  force vector/.style={
    -{Latex[length=3.2mm]},
    line width=1.45pt,
    draw=visualred
  },
  field vector/.style={
    -{Latex[length=2.8mm]},
    line width=1.1pt,
    draw=visualgreen
  },
  construction line/.style={
    dashed,
    line width=0.75pt,
    draw=visualgrid
  },
  trajectory/.style={
    line width=1.5pt,
    draw=visualblue
  },
  visual label/.style={
    font=\small,
    text=visualink,
    fill=white,
    inner sep=1.5pt
  },
  concept node/.style={
    rounded corners=2pt,
    draw=visualblue,
    fill=softblue,
    line width=0.8pt,
    align=center,
    inner sep=5pt,
    font=\small
  }
}

% Figure caption system ---------------------------------------------------------
\captionsetup{
  font=small,
  labelfont={bf,color=visualblue},
  textfont={color=visualink},
  justification=raggedright,
  singlelinecheck=false,
  skip=6pt
}

\newcommand{\figsource}[1]{%
  \par\smallskip
  {\footnotesize\color{gray}\textit{Source/Construction:} #1}
}

\newenvironment{bookfigure}[2][tbp]{%
  \def\bookfigurecaption{#2}%
  \begin{figure}[#1]
  \centering
  \begin{adjustbox}{max width=\textwidth}
}{%
  \end{adjustbox}
  \caption{\bookfigurecaption}
  \end{figure}
}

% Compact visual callout --------------------------------------------------------
\newtcolorbox{visualguidebox}[1][How to read the figure]{
  title=#1,
  colback=white,
  colframe=visualcyan,
  fonttitle=\bfseries,
  borderline west={2pt}{0pt}{visualcyan}
}

% Chapter-opening artwork -------------------------------------------------------
\newcommand{\chapteropeningart}[2]{%
\begin{center}
\begin{tikzpicture}[x=1cm,y=1cm]
  \fill[softblue] (0,0) rectangle (12.6,2.65);
  \draw[visualblue,line width=1pt] (0,0) rectangle (12.6,2.65);
  \draw[visualgrid,line width=0.5pt] (0.45,0.35) grid[xstep=0.55,ystep=0.55] (12.1,2.3);
  \node[anchor=west,font=\bfseries\Large,text=visualblue] at (0.7,1.78) {#1};
  \node[anchor=west,align=left,text width=7.2cm,font=\small,text=visualink]
    at (0.72,0.88) {#2};
  \draw[physics axis] (9.1,0.62) -- (11.85,0.62) node[right] {$x$};
  \draw[physics axis] (9.1,0.62) -- (9.1,2.18) node[above] {$y$};
  \draw[physics vector] (9.1,0.62) -- (11.08,1.73)
    node[above right,visual label] {$\mathbf q$};
  \fill[visualgold] (9.1,0.62) circle (1.7pt);
\end{tikzpicture}
\end{center}
\vspace{0.4cm}
}

% Future-ready visual templates -------------------------------------------------
\newcommand{\wavefunctionplot}{%
\begin{tikzpicture}
  \begin{axis}[
    width=10.5cm,
    height=5.2cm,
    axis lines=middle,
    xlabel={$x$},
    ylabel={$\psi(x)$},
    ticks=none,
    domain=-3.3:3.3,
    samples=180,
    ymin=-0.75,
    ymax=1.15
  ]
    \addplot[visualblue,very thick] {exp(-x^2/2)};
    \addplot[visualcyan,thick,dashed] {0.72*x*exp(-x^2/2)};
  \end{axis}
\end{tikzpicture}
}

\newcommand{\quantumhalltemplate}{%
\begin{tikzpicture}[x=1cm,y=1cm]
  \fill[softblue] (0,0) rectangle (9.8,4.3);
  \draw[visualblue,line width=0.9pt] (0,0) rectangle (9.8,4.3);
  \foreach \x in {0.7,1.7,...,9.2}
    \foreach \y in {0.6,1.6,...,3.8}
      \draw[visualcyan!55,line width=0.5pt] (\x,\y) circle (0.16);
  \draw[force vector] (0.55,3.75) -- (9.15,3.75)
    node[midway,above,visual label] {upper chiral edge state};
  \draw[force vector] (9.15,0.55) -- (0.55,0.55)
    node[midway,below,visual label] {lower chiral edge state};
  \node[font=\small] at (4.9,2.15) {2DEG bulk: Landau-level states};
  \draw[physics vector] (8.85,1.5) -- (8.85,2.85)
    node[above] {$\mathbf B$};
\end{tikzpicture}
}

\title{\Huge\bfseries Theory of Quantum Mechanics\\[0.55em]
\Large From Mathematical Foundations to the Quantum Hall Effect}
\author{}
\date{PR-BOOK-04 Visual System\\\today}

\begin{document}
	
	
%%%%%%%%%%%%%%%%%%%%%%%%%%%%%%%%%%%%%%%%%%%%%%%%%%%%%%%%%%%%%%%%%%%%%%%%%%%%%%%%
% Front Matter
%%%%%%%%%%%%%%%%%%%%%%%%%%%%%%%%%%%%%%%%%%%%%%%%%%%%%%%%%%%%%%%%%%%%%%%%%%%%%%%%

\frontmatter

\maketitle

\clearpage

\pdfbookmark[0]{Table of Contents}{toc}
\tableofcontents

\clearpage

\chapter*{Preface}
\addcontentsline{toc}{chapter}{Preface}

This book develops the mathematical and physical ideas needed to understand
modern quantum mechanics and the Quantum Hall Effect. It begins with vectors,
motion, and calculus, then advances through classical mechanics,
electromagnetism, analytical mechanics, and quantum theory.

The presentation emphasizes geometric intuition, careful derivation, and
explicit links between classical and quantum descriptions. Each chapter uses a
consistent structure: learning objectives, conceptual development, physical
interpretation, worked examples, common mistakes, a summary, and a preview of
the next chapter.

\clearpage

%%%%%%%%%%%%%%%%%%%%%%%%%%%%%%%%%%%%%%%%%%%%%%%%%%%%%%%%%%%%%%%%%%%%%%%%%%%%%%%%
% Main Matter
%%%%%%%%%%%%%%%%%%%%%%%%%%%%%%%%%%%%%%%%%%%%%%%%%%%%%%%%%%%%%%%%%%%%%%%%%%%%%%%%
	
%==============================================================================
% PR-BOOK-01 : Front Matter
%==============================================================================

\frontmatter

%%%%%%%%%%%%%%%%%%%%%%%%%%%%%%%%%%%%%%%%%%%%%%%%%%%%%%%%%%%%%%%%%%%%%%%%%%%%%%%
% Half Title
%%%%%%%%%%%%%%%%%%%%%%%%%%%%%%%%%%%%%%%%%%%%%%%%%%%%%%%%%%%%%%%%%%%%%%%%%%%%%%%

\begin{titlepage}
	\centering
	\vspace*{\fill}
	
	{\Huge\bfseries Theory of Quantum Mechanics\par}
	
	\vspace{1cm}
	
	{\Large Volume I}
	
	\vspace*{\fill}
\end{titlepage}

%%%%%%%%%%%%%%%%%%%%%%%%%%%%%%%%%%%%%%%%%%%%%%%%%%%%%%%%%%%%%%%%%%%%%%%%%%%%%%%
% Full Title
%%%%%%%%%%%%%%%%%%%%%%%%%%%%%%%%%%%%%%%%%%%%%%%%%%%%%%%%%%%%%%%%%%%%%%%%%%%%%%%

\begin{titlepage}
	
	\centering
	
	\vspace*{2cm}
	
	{\Huge\bfseries Theory of Quantum Mechanics\par}
	
	\vspace{0.8cm}
	
	{\Large From Mathematical Foundations\\[0.2cm]
		to the Quantum Hall Effect\par}
	
	\vspace{1cm}
	
	{\Large Volume I\\
		Mathematical Foundations\par}
	
	\vspace{2cm}
	
	Version 0.9 Draft
	
	\vfill
	
	Author
	
	\vspace{0.5cm}
	
	\today
	
\end{titlepage}

%%%%%%%%%%%%%%%%%%%%%%%%%%%%%%%%%%%%%%%%%%%%%%%%%%%%%%%%%%%%%%%%%%%%%%%%%%%%%%%
% Copyright
%%%%%%%%%%%%%%%%%%%%%%%%%%%%%%%%%%%%%%%%%%%%%%%%%%%%%%%%%%%%%%%%%%%%%%%%%%%%%%%

\cleardoublepage

\thispagestyle{empty}

\section*{Copyright}

Copyright \copyright\ 2026.

All rights reserved.

This manuscript is a working draft intended for educational and research
purposes.

Future editions may contain additional chapters, figures, exercises and
historical notes.

\vfill

ISBN: XXX-X-XXXX-XXXX-X

Publisher: To be assigned

%%%%%%%%%%%%%%%%%%%%%%%%%%%%%%%%%%%%%%%%%%%%%%%%%%%%%%%%%%%%%%%%%%%%%%%%%%%%%%%
% Dedication
%%%%%%%%%%%%%%%%%%%%%%%%%%%%%%%%%%%%%%%%%%%%%%%%%%%%%%%%%%%%%%%%%%%%%%%%%%%%%%%

\cleardoublepage

\thispagestyle{empty}

\vspace*{\fill}

\begin{center}
	\itshape
	
	Dedicated to everyone who wants to understand\\
	Nature through mathematics.
	
\end{center}

\vspace*{\fill}

%%%%%%%%%%%%%%%%%%%%%%%%%%%%%%%%%%%%%%%%%%%%%%%%%%%%%%%%%%%%%%%%%%%%%%%%%%%%%%%
% Preface
%%%%%%%%%%%%%%%%%%%%%%%%%%%%%%%%%%%%%%%%%%%%%%%%%%%%%%%%%%%%%%%%%%%%%%%%%%%%%%%

\chapter*{Preface}
\addcontentsline{toc}{chapter}{Preface}

Quantum mechanics is one of the greatest intellectual achievements in science.
Rather than presenting the subject as a collection of formulas, this book
builds the theory from first principles.

The journey begins with vectors, geometry and calculus before progressing to

\begin{itemize}
	\item classical mechanics,
	\item electromagnetism,
	\item analytical mechanics,
	\item quantum mechanics,
	\item Landau levels,
	\item the Quantum Hall Effect,
	\item modern topological physics.
\end{itemize}

Every chapter follows the same educational structure:

\begin{enumerate}
	\item Learning Objectives
	\item Conceptual Motivation
	\item Definitions
	\item Theory
	\item Worked Examples
	\item Physical Interpretation
	\item Historical Notes
	\item Common Mistakes
	\item Summary
	\item Exercises
	\item Looking Ahead
\end{enumerate}

%%%%%%%%%%%%%%%%%%%%%%%%%%%%%%%%%%%%%%%%%%%%%%%%%%%%%%%%%%%%%%%%%%%%%%%%%%%%%%%
% How to Use This Book
%%%%%%%%%%%%%%%%%%%%%%%%%%%%%%%%%%%%%%%%%%%%%%%%%%%%%%%%%%%%%%%%%%%%%%%%%%%%%%%

\chapter*{How to Use This Book}
\addcontentsline{toc}{chapter}{How to Use This Book}

Different readers may use this text in different ways.

\begin{description}
	\item[First reading:] Focus on intuition and worked examples.
	\item[Second reading:] Study the mathematical derivations.
	\item[Reference:] Use the notation tables, summaries and index to locate concepts quickly.
\end{description}

%%%%%%%%%%%%%%%%%%%%%%%%%%%%%%%%%%%%%%%%%%%%%%%%%%%%%%%%%%%%%%%%%%%%%%%%%%%%%%%
% Frequently Used Symbols
%%%%%%%%%%%%%%%%%%%%%%%%%%%%%%%%%%%%%%%%%%%%%%%%%%%%%%%%%%%%%%%%%%%%%%%%%%%%%%%

\chapter*{Frequently Used Symbols}
\addcontentsline{toc}{chapter}{Frequently Used Symbols}

\begin{tabular}{lll}
	\textbf{Symbol} & \textbf{Meaning} & \textbf{SI Unit}\\
	\hline
	$\mathbf r$ & Position vector & m\\
	$\mathbf v$ & Velocity & m/s\\
	$\mathbf a$ & Acceleration & m/s$^2$\\
	$\mathbf F$ & Force & N\\
	$\mathbf p$ & Momentum & kg\,m/s\\
	$\nabla$ & Nabla operator & --\\
	$\psi$ & Wavefunction & --\\
\end{tabular}

%%%%%%%%%%%%%%%%%%%%%%%%%%%%%%%%%%%%%%%%%%%%%%%%%%%%%%%%%%%%%%%%%%%%%%%%%%%%%%%
% Physical Constants
%%%%%%%%%%%%%%%%%%%%%%%%%%%%%%%%%%%%%%%%%%%%%%%%%%%%%%%%%%%%%%%%%%%%%%%%%%%%%%%

\chapter*{Fundamental Physical Constants}
\addcontentsline{toc}{chapter}{Fundamental Physical Constants}

\begin{tabular}{lll}
	Constant & Symbol & Value\\
	\hline
	Speed of light & $c$ & $2.99792458\times10^8$ m/s\\
	Planck constant & $h$ & $6.62607015\times10^{-34}$ J\,s\\
	Reduced Planck constant & $\hbar$ & $1.054571817\times10^{-34}$ J\,s\\
	Elementary charge & $e$ & $1.602176634\times10^{-19}$ C\\
	Electron mass & $m_e$ & $9.1093837015\times10^{-31}$ kg\\
	Vacuum permittivity & $\varepsilon_0$ & $8.8541878128\times10^{-12}$ F/m\\
\end{tabular}

%%%%%%%%%%%%%%%%%%%%%%%%%%%%%%%%%%%%%%%%%%%%%%%%%%%%%%%%%%%%%%%%%%%%%%%%%%%%%%%
% Roadmap
%%%%%%%%%%%%%%%%%%%%%%%%%%%%%%%%%%%%%%%%%%%%%%%%%%%%%%%%%%%%%%%%%%%%%%%%%%%%%%%

\chapter*{Roadmap of the Book}
\addcontentsline{toc}{chapter}{Roadmap}

Part I --- Mathematical Foundations

Part II --- Classical Mechanics

Part III --- Electromagnetism

Part IV --- Analytical Mechanics

Part V --- Quantum Mechanics

Part VI --- Quantum Hall Effect

Part VII --- Topology and Modern Quantum Matter

\cleardoublepage
	
	
	
%%%%%%%%%%%%%%%%%%%%%%%%%%%%%%%%%%%%%%%%%%%%%%%%%%%%%%%%%%%%%%%%%%%%%%%%%%%%%%%%
\mainmatter

\chapter{Vectors and Coordinate Systems}
\chapteropeningart
{Geometry becomes physics}
{Vectors encode magnitude and direction. Coordinate systems provide numerical
components, while the underlying geometric object remains independent of the
chosen axes.}
%%%%%%%%%%%%%%%%%%%%%%%%%%%%%%%%%%%%%%%%%%%%%%%%%%%%%%%%%%%%%%%%%%%%%%%%%%%%%%%%

\chapterquote
{``The laws of physics take the same form in every coordinate system.
Coordinates are chosen by us; nature does not know them.''}
{Adapted from Albert Einstein}

\begin{roadmapbox}
Scalars $\longrightarrow$ vectors $\longrightarrow$ coordinates
$\longrightarrow$ vector operations $\longrightarrow$ physical fields.
\end{roadmapbox}

\learningobjectives


\begin{itemize}
	\item understand why vectors are necessary in physics;
	\item distinguish scalar and vector quantities;
	\item work with Cartesian coordinate systems;
	\item calculate vector magnitudes;
	\item perform vector addition and subtraction;
	\item calculate dot and cross products;
	\item understand the right-hand rule;
	\item interpret vectors geometrically;
	\item understand basis vectors;
	\item transform between coordinate systems;
	\item appreciate why vectors appear naturally in Maxwell's equations,
	Newton's laws and Schrödinger's equation.
\end{itemize}

\newpage

%%%%%%%%%%%%%%%%%%%%%%%%%%%%%%%%%%%%%%%%%%%%%%%%%%%%%%

\section{Why Do We Need Vectors?}

Suppose someone tells you

\[
v=20\ \mathrm{m/s}.
\]

Has the motion been completely described?

No.

The object could move

east,

west,

north,

south,

up,

or any other direction.

Speed alone is insufficient.

Now suppose someone writes

\[
\mathbf v=
20\,\mathrm{m/s}
\ \hat{\mathbf x}.
\]

Now the motion is completely specified.

The velocity has

\begin{itemize}
	\item magnitude,
	\item direction.
\end{itemize}

This simple observation motivates one of the most important concepts in
physics: the vector.

%%%%%%%%%%%%%%%%%%%%%%%%%%%%%%%%%%%%%%%%%%%%%%%%%%%%%%

\section{Scalars}

A scalar possesses only magnitude.

Examples include

\begin{itemize}
	\item mass,
	\item temperature,
	\item electric potential,
	\item pressure,
	\item energy,
	\item probability density.
\end{itemize}

Scalars do not possess direction.

Examples

\[
T=300\ {\rm K},
\]

\[
m=9.11\times10^{-31}\ {\rm kg}.
\]

%%%%%%%%%%%%%%%%%%%%%%%%%%%%%%%%%%%%%%%%%%%%%%%%%%%%%%

\section{Vectors}

Vectors possess both

\begin{itemize}
	\item magnitude,
	\item direction.
\end{itemize}

Examples include

\begin{itemize}
	\item displacement,
	\item velocity,
	\item acceleration,
	\item force,
	\item momentum,
	\item electric field,
	\item magnetic field,
	\item probability current.
\end{itemize}

Vectors are usually written in boldface,

\[
\mathbf A,
\]

or with an arrow,

\[
\vec A.
\]

%%%%%%%%%%%%%%%%%%%%%%%%%%%%%%%%%%%%%%%%%%%%%%%%%%%%%%


\begin{bookfigure}{A vector resolved into Cartesian components. The same
geometric vector can be described by different components after a change of
coordinates.}
\begin{tikzpicture}[scale=1.05]
  \draw[physics axis] (-0.2,0) -- (5.6,0) node[right] {$x$};
  \draw[physics axis] (0,-0.2) -- (0,4.2) node[above] {$y$};
  \draw[physics vector] (0,0) -- (4.4,3.15)
    node[above right,visual label] {$\mathbf A$};
  \draw[construction line] (4.4,0) -- (4.4,3.15);
  \draw[construction line] (0,3.15) -- (4.4,3.15);
  \draw[physics vector secondary] (0,-0.25) -- (4.4,-0.25)
    node[midway,below] {$A_x$};
  \draw[physics vector secondary] (-0.25,0) -- (-0.25,3.15)
    node[midway,left] {$A_y$};
  \draw[visualgold,line width=0.9pt,-{Latex[length=1.5mm]}]
    (0.9,0) arc[start angle=0,end angle=35.6,radius=0.9];
  \node[text=visualgold,font=\small] at (1.12,0.36) {$\theta$};
\end{tikzpicture}
\end{bookfigure}

\begin{visualguidebox}
The diagonal arrow is the vector itself. The horizontal and vertical arrows are
its components relative to the selected Cartesian basis.
\end{visualguidebox}


\section{Space}

Physics takes place in space.

To describe position we choose coordinates.

Nature itself has no coordinate axes.

Coordinates are introduced only to help us describe observations.

The most commonly used coordinate system is Cartesian coordinates.

A point is represented by

\[
(x,y,z).
\]

These numbers describe

\begin{itemize}
	\item distance along the x-axis,
	\item distance along the y-axis,
	\item distance along the z-axis.
\end{itemize}

%%%%%%%%%%%%%%%%%%%%%%%%%%%%%%%%%%%%%%%%%%%%%%%%%%%%%%

\section{Cartesian Coordinate System}

Define three mutually perpendicular axes.

\[
x,
\qquad
y,
\qquad
z.
\]

Their associated unit vectors are

\[
\hat{\mathbf x},
\qquad
\hat{\mathbf y},
\qquad
\hat{\mathbf z}.
\]

Each has unit length

\[
|\hat{\mathbf x}|=
|\hat{\mathbf y}|=
|\hat{\mathbf z}|=1.
\]

They satisfy

\[
\hat{\mathbf x}\cdot\hat{\mathbf y}=0,
\]

etc.

Therefore they are orthonormal.

%%%%%%%%%%%%%%%%%%%%%%%%%%%%%%%%%%%%%%%%%%%%%%%%%%%%%%

\section{Position Vector}

A point

\[
(x,y,z)
\]

can also be written as

\[
\boxed{
	\mathbf r
	=
	x\hat{\mathbf x}
	+
	y\hat{\mathbf y}
	+
	z\hat{\mathbf z}
}
\]

This is called the position vector.

Notice that the coordinates depend on the chosen axes.

The physical point does not.

%%%%%%%%%%%%%%%%%%%%%%%%%%%%%%%%%%%%%%%%%%%%%%%%%%%%%%

\section{Magnitude of a Vector}

Consider

\[
\mathbf A=
A_x\hat x
+
A_y\hat y
+
A_z\hat z.
\]

The magnitude is

\[
|\mathbf A|
=
\sqrt{
	A_x^2
	+
	A_y^2
	+
	A_z^2
}.
\]

This is simply the three-dimensional Pythagorean theorem.

%%%%%%%%%%%%%%%%%%%%%%%%%%%%%%%%%%%%%%%%%%%%%%%%%%%%%%

\section{Unit Vectors}

A unit vector has length one.

Every vector can be written as

\[
\mathbf A
=
|\mathbf A|
\hat{\mathbf A}.
\]

Therefore

\[
\hat{\mathbf A}
=
\frac{\mathbf A}{|\mathbf A|}.
\]

%%%%%%%%%%%%%%%%%%%%%%%%%%%%%%%%%%%%%%%%%%%%%%%%%%%%%%

\section{Vector Addition}

Vectors add component by component.

\[
\mathbf C
=
\mathbf A+\mathbf B.
\]

Therefore

\[
C_x=A_x+B_x,
\]

etc.

%%%%%%%%%%%%%%%%%%%%%%%%%%%%%%%%%%%%%%%%%%%%%%%%%%%%%%

\section{Dot Product}

The scalar product is

\[
\boxed{
	\mathbf A\cdot\mathbf B
	=
	A_xB_x
	+
	A_yB_y
	+
	A_zB_z
}
\]

or

\[
\boxed{
	\mathbf A\cdot\mathbf B
	=
	|\mathbf A|
	|\mathbf B|
	\cos\theta.
}
\]

Applications include

\begin{itemize}
	\item work,
	\item projections,
	\item orthogonality.
\end{itemize}

%%%%%%%%%%%%%%%%%%%%%%%%%%%%%%%%%%%%%%%%%%%%%%%%%%%%%%

\section{Cross Product}

The vector product is

\[
\boxed{
	\mathbf A\times\mathbf B
}
\]

whose magnitude equals

\[
|\mathbf A\times\mathbf B|
=
AB\sin\theta.
\]

Its direction follows the right-hand rule.

The result is perpendicular to both vectors.

%%%%%%%%%%%%%%%%%%%%%%%%%%%%%%%%%%%%%%%%%%%%%%%%%%%%%%

\section{The Right-Hand Rule}

The right-hand rule appears throughout electromagnetism.

Examples include

\[
\mathbf F=q\mathbf v\times\mathbf B,
\]

\[
\mathbf B=\nabla\times\mathbf A,
\]

and angular momentum

\[
\mathbf L
=
\mathbf r\times\mathbf p.
\]

Later, this rule will determine the direction of the Lorentz force in the Hall effect.

%%%%%%%%%%%%%%%%%%%%%%%%%%%%%%%%%%%%%%%%%%%%%%%%%%%%%%

% PR-BOOK-07.2: canonical modular section
% PR-BOOK-07.2
% Canonical Chapter 1 section: Vector Projection and Decomposition
% This file is included by main.tex after the dot and cross product material.

\section{Vector Projection and Decomposition}
\label{sec:vector-projection-decomposition}

One of the most powerful techniques in physics is to decompose a complicated
quantity into simpler parts. Instead of analysing an entire vector at once, we
often separate it into components that have clear physical meaning.

Consider, for example, a block resting on an inclined plane. The gravitational
force acts vertically downward, yet only part of this force pulls the block
down the slope. Another part presses the block against the surface. By resolving
the gravitational force into two perpendicular components, the problem becomes
much easier to analyse.

This idea appears throughout science and engineering. In mechanics we resolve
forces and velocities; in electromagnetism we resolve electric and magnetic
fields; and in quantum mechanics we project state vectors onto basis states.
All of these applications rely on the mathematical concept of projection.

\subsection{Motivation}

Suppose that a force $\mathbf F$ acts on an object, but that we are interested
only in the part of the force acting along a direction represented by a non-zero
vector $\mathbf a$. The relevant question is not ``What is the total force?''
but rather ``How much of this force acts in the chosen direction?''

The answer is provided by the projection of one vector onto another.

\subsection{Scalar Projection}

Let $\mathbf A$ and $\mathbf B$ be non-zero vectors separated by an angle
$\theta$. The \emph{scalar projection}, also called the scalar component, of
$\mathbf A$ onto $\mathbf B$ is

\[
\boxed{
\operatorname{comp}_{\mathbf B}\mathbf A
=
\frac{\mathbf A\cdot\mathbf B}{\lVert\mathbf B\rVert}
=
\lVert\mathbf A\rVert\cos\theta
}
\]

The result is a scalar quantity with the same physical units as the original
vector. Geometrically, it represents the signed length of the shadow cast by
$\mathbf A$ onto the direction of $\mathbf B$. The sign distinguishes alignment
from anti-alignment.

\begin{figure}[htbp]
\centering
\begingroup
% Local compatibility aliases for the reusable Chapter 1 illustration library.
% They are scoped to this figure and do not change the global visual system.
\colorlet{PhysicsBlue}{visualblue}
\colorlet{MotionOrange}{visualcyan}
\colorlet{ForceRed}{visualred}
\colorlet{FieldGreen}{visualgreen}
\colorlet{ConstructionGray}{visualgrid}
\colorlet{SoftBlue}{softblue}
\tikzset{
  axis/.style={physics axis},
  vector/.style={physics vector},
  vector secondary/.style={physics vector secondary},
  force/.style={force vector},
  field/.style={field vector},
  construction/.style={construction line},
  label/.style={visual label}
}
\begin{tikzpicture}[scale=1.05]
  \coordinate (O) at (0,0);
  \coordinate (B) at (5.3,1.2);
  \coordinate (A) at (3.4,3.65);
  \draw[vector secondary] (O) -- (B) node[right,label] {$\mathbf b$};
  \draw[vector] (O) -- (A) node[above,label] {$\mathbf a$};
  \coordinate (P) at ($(O)!(A)!(B)$);
  \draw[construction] (A) -- (P);
  \draw[force] (O) -- (P)
    node[midway,below,label] {$\operatorname{proj}_{\mathbf b}\mathbf a$};
  \draw[PhysicsBlue,line width=1pt]
    (1.15,.26) arc[start angle=12.8,end angle=47.1,radius=1.18];
  \node[label] at (1.32,.62) {$\theta$};
  \node[concept node,text width=4.9cm,anchor=west] at (6.0,2.1) {
    Scalar projection:
    $\displaystyle \frac{\mathbf a\cdot\mathbf b}{\|\mathbf b\|}$\\[1mm]
    Vector projection:
    $\displaystyle\frac{\mathbf a\cdot\mathbf b}{\|\mathbf b\|^2}\mathbf b$
  };
\end{tikzpicture}

\endgroup
\caption{Scalar and vector projection of $\mathbf A$ onto the direction of
$\mathbf B$. The dashed segment is perpendicular to $\mathbf B$.}
\label{fig:vector-projection}
\end{figure}

\subsection{Vector Projection}

Frequently we require not only the length of the projection but also its
direction. The \emph{vector projection} of $\mathbf A$ onto $\mathbf B$ is

\[
\boxed{
\operatorname{proj}_{\mathbf B}\mathbf A
=
\left(
\frac{\mathbf A\cdot\mathbf B}{\mathbf B\cdot\mathbf B}
\right)\mathbf B
}
\]

Unlike the scalar projection, this expression is itself a vector. It is parallel
to $\mathbf B$, and its signed magnitude is determined by the scalar projection.

\subsection{Derivation of the Projection Formula}

The projection formula follows directly from the dot product. First determine
the signed length of the projection:

\[
\lVert\mathbf A\rVert\cos\theta
=
\frac{\mathbf A\cdot\mathbf B}{\lVert\mathbf B\rVert}.
\]

The unit vector in the direction of $\mathbf B$ is

\[
\widehat{\mathbf B}
=
\frac{\mathbf B}{\lVert\mathbf B\rVert}.
\]

Multiplying the projected length by this unit vector gives

\[
\operatorname{proj}_{\mathbf B}\mathbf A
=
\frac{\mathbf A\cdot\mathbf B}{\lVert\mathbf B\rVert}
\frac{\mathbf B}{\lVert\mathbf B\rVert}.
\]

Since $\lVert\mathbf B\rVert^2=\mathbf B\cdot\mathbf B$, this becomes

\[
\boxed{
\operatorname{proj}_{\mathbf B}\mathbf A
=
\frac{\mathbf A\cdot\mathbf B}{\mathbf B\cdot\mathbf B}\mathbf B.
}
\]

Thus, projection is not an unrelated formula to be memorized; it is a direct
consequence of the geometric meaning of the dot product.

\subsection{Orthogonal Decomposition}

Every vector $\mathbf A$ may be decomposed uniquely into two parts relative to
a chosen non-zero direction $\mathbf B$:

\[
\boxed{
\mathbf A=\mathbf A_{\parallel}+\mathbf A_{\perp}.
}
\]

The component parallel to $\mathbf B$ is

\[
\boxed{
\mathbf A_{\parallel}
=
\operatorname{proj}_{\mathbf B}\mathbf A,
}
\]

while the perpendicular component is obtained by subtraction:

\[
\boxed{
\mathbf A_{\perp}
=
\mathbf A-\mathbf A_{\parallel}.
}
\]

By construction,

\[
\mathbf A_{\perp}\cdot\mathbf B=0.
\]

This orthogonality can be verified directly:

\[
\begin{aligned}
\mathbf A_{\perp}\cdot\mathbf B
&=
\left(
\mathbf A
-
\frac{\mathbf A\cdot\mathbf B}{\mathbf B\cdot\mathbf B}\mathbf B
\right)\cdot\mathbf B
\\
&=
\mathbf A\cdot\mathbf B
-
\frac{\mathbf A\cdot\mathbf B}{\mathbf B\cdot\mathbf B}
(\mathbf B\cdot\mathbf B)
\\
&=0.
\end{aligned}
\]

\begin{figure}[htbp]
\centering
\begingroup
% Local compatibility aliases for the reusable Chapter 1 illustration library.
% They are scoped to this figure and do not change the global visual system.
\colorlet{PhysicsBlue}{visualblue}
\colorlet{MotionOrange}{visualcyan}
\colorlet{ForceRed}{visualred}
\colorlet{FieldGreen}{visualgreen}
\colorlet{ConstructionGray}{visualgrid}
\colorlet{SoftBlue}{softblue}
\tikzset{
  axis/.style={physics axis},
  vector/.style={physics vector},
  vector secondary/.style={physics vector secondary},
  force/.style={force vector},
  field/.style={field vector},
  construction/.style={construction line},
  label/.style={visual label}
}
% Orthogonal decomposition of a vector into Cartesian components
\begin{tikzpicture}[scale=1.0]
  \draw[axis] (-.45,0)--(5.8,0) node[right] {$x$};
  \draw[axis] (0,-.45)--(0,4.5) node[above] {$y$};
  \coordinate (P) at (4.6,3.25);
  \draw[vector] (0,0)--(P) node[midway,above left,label] {$\mathbf A$};
  \draw[vector secondary] (0,0)--(4.6,0) node[midway,below,label] {$A_x\hat{\mathbf i}$};
  \draw[field] (4.6,0)--(P) node[midway,right,label] {$A_y\hat{\mathbf j}$};
  \draw[construction] (0,3.25)--(P);
  \draw[PhysicsBlue!65,line width=.8pt] (.9,0) arc[start angle=0,end angle=35.25,radius=.9];
  \node[label] at (1.15,.38) {$\theta$};
  \node[concept node,text width=4.4cm] at (7.9,2.0) {$A_x=A\cos\theta$\\$A_y=A\sin\theta$\\[1mm]$\mathbf A=A_x\hat{\mathbf i}+A_y\hat{\mathbf j}$};
\end{tikzpicture}

\endgroup
\caption{A vector decomposed into two mutually perpendicular components.}
\label{fig:orthogonal-decomposition}
\end{figure}

\subsection{Application: Force on an Inclined Plane}

Consider a block of mass $m$ resting on a plane inclined at an angle $\alpha$
above the horizontal. The gravitational force is

\[
\mathbf F_g=m\mathbf g.
\]

Rather than analysing this force directly, we resolve it into components parallel
and perpendicular to the plane. Their magnitudes are

\[
\boxed{F_{\parallel}=mg\sin\alpha}
\]

and

\[
\boxed{F_{\perp}=mg\cos\alpha.}
\]

The parallel component tends to accelerate the block down the incline. The
perpendicular component determines the normal force exerted by the surface,
provided that no other force has a component normal to the plane.

This decomposition converts a two-dimensional force problem into two coupled
but individually one-dimensional equations.

\subsection{Worked Example}

Let

\[
\mathbf A=(3,4),
\qquad
\mathbf B=(5,0).
\]

Determine the projection of $\mathbf A$ onto $\mathbf B$ and the component of
$\mathbf A$ perpendicular to $\mathbf B$.

First compute the dot products:

\[
\mathbf A\cdot\mathbf B
=
3(5)+4(0)
=
15,
\]

and

\[
\mathbf B\cdot\mathbf B
=
5^2+0^2
=
25.
\]

Therefore,

\[
\operatorname{proj}_{\mathbf B}\mathbf A
=
\frac{15}{25}(5,0)
=
(3,0).
\]

The perpendicular component is

\[
\mathbf A_{\perp}
=
(3,4)-(3,0)
=
(0,4).
\]

Finally,

\[
(3,0)\cdot(0,4)=0,
\]

which confirms that the two components are orthogonal.

\subsection{Physical Interpretation}

Projection isolates the part of a vector that is physically relevant to a
chosen direction. Important examples include:

\begin{itemize}
\item the component of velocity directed toward an observer in the Doppler effect;
\item the component of force along a displacement in mechanical work;
\item the electric field normal to a conducting surface;
\item the magnetic field perpendicular to the motion of a charged particle;
\item the amplitude of a quantum state projected onto a measurement basis.
\end{itemize}

Although these examples belong to different branches of physics, they employ
the same mathematical operation.

\subsection{Historical Note}

The systematic use of orthogonal decomposition developed during the nineteenth
century through the work of Hermann Grassmann and later Josiah Willard Gibbs.
In the twentieth century, David Hilbert generalized these ideas to
infinite-dimensional inner-product spaces. In quantum mechanics, measurement
is represented mathematically through projections onto subspaces associated
with possible outcomes.

\subsection{Common Mistakes}

The scalar projection and vector projection are different quantities. The first
is a signed number; the second is a vector. Projection also depends on the
direction onto which the vector is projected, so in general

\[
\operatorname{proj}_{\mathbf B}\mathbf A
\neq
\operatorname{proj}_{\mathbf A}\mathbf B.
\]

Finally, the perpendicular component should be calculated by subtracting the
parallel projection from the original vector:

\[
\mathbf A_{\perp}
=
\mathbf A-\operatorname{proj}_{\mathbf B}\mathbf A.
\]

\subsection{Summary}

The dot product measures alignment between vectors. A scalar projection extracts
the signed magnitude of one vector along another, while a vector projection
includes both magnitude and direction. Every vector can be decomposed uniquely
into parallel and perpendicular parts relative to a non-zero reference
direction. This decomposition simplifies physical problems by separating
independent directions.

These ideas will be used repeatedly in Newtonian mechanics, electromagnetism,
wave mechanics, and quantum theory.

\subsection*{Looking Ahead}

The next stage develops coordinate transformations, including rotations,
reflections, translations, and changes of basis. Projection and decomposition
will provide the conceptual bridge between geometric vectors and their
coordinate representations.


\subsection*{Projection Operators and Quantum Preview}
Projection can be viewed as an operator $P$ acting on vectors. A projection operator satisfies $P^2=P$, meaning that applying the same projection twice has the same effect as applying it once. In linear algebra and quantum mechanics, projection operators play a central role in describing measurements.

Repeated projections can be used to construct orthogonal bases through the Gram--Schmidt orthogonalization procedure. In quantum mechanics, every measurement may be interpreted as projecting a state vector onto an orthonormal basis. The probability of observing a particular outcome is given by the squared magnitude of the corresponding projection, providing a direct bridge from elementary vector algebra to Hilbert spaces.


\section{Preview of Quantum Mechanics}

Although wavefunctions are complex-valued rather than ordinary vectors,
many quantum quantities are vectors:

\[
\nabla,
\qquad
\mathbf p,
\qquad
\mathbf A,
\qquad
\mathbf j.
\]

Understanding vectors is therefore essential before studying the Schrödinger equation.

%%%%%%%%%%%%%%%%%%%%%%%%%%%%%%%%%%%%%%%%%%%%%%%%%%%%%%

\section{Summary}

This chapter introduced the mathematical language used throughout the remainder of the book.

The reader should now understand

\begin{itemize}
	\item scalar quantities,
	\item vectors,
	\item coordinate systems,
	\item basis vectors,
	\item vector magnitude,
	\item vector addition,
	\item dot product,
	\item cross product,
	\item the right-hand rule.
\end{itemize}

These concepts will be used immediately in the next chapter, where we introduce derivatives, vector calculus, and the differential operator

\[
\nabla.
\]


%%%%%%%%%%%%%%%%%%%%%%%%%%%%%%%%%%%%%%%%%%%%%%%%%%%%%%%%%%%%%%%%%%%%%%%%%%%%%%%%


\begin{bookfigure}{Geometric meanings of the dot and cross products.}
\begin{tikzpicture}[x=1cm,y=1cm]
  % Dot product panel
  \node[font=\bfseries,text=visualblue] at (2.3,3.3) {Dot product};
  \draw[physics vector] (0.4,0.6) -- (3.8,0.6) node[right] {$\mathbf A$};
  \draw[physics vector secondary] (0.4,0.6) -- (3.05,2.55) node[above] {$\mathbf B$};
  \draw[construction line] (3.05,2.55) -- (3.05,0.6);
  \draw[line width=2.2pt,draw=visualgold] (0.4,0.35) -- (3.05,0.35);
  \node[font=\small] at (1.75,0.05) {projection of $\mathbf B$ on $\mathbf A$};

  % Cross product panel
  \node[font=\bfseries,text=visualblue] at (7.7,3.3) {Cross product};
  \draw[physics vector] (5.4,0.65) -- (9.2,0.65) node[right] {$\mathbf A$};
  \draw[physics vector secondary] (5.4,0.65) -- (7.7,2.65) node[above] {$\mathbf B$};
  \fill[visualcyan!18] (5.4,0.65) -- (9.2,0.65) -- (11.5,2.65) -- (7.7,2.65) -- cycle;
  \draw[visualcyan,line width=0.8pt] (5.4,0.65) -- (9.2,0.65) -- (11.5,2.65) -- (7.7,2.65) -- cycle;
  \draw[force vector] (8.1,1.55) -- (8.1,3.05)
    node[above] {$\mathbf A\times\mathbf B$};
  \node[font=\small] at (8.4,0.15) {magnitude = parallelogram area};
\end{tikzpicture}
\end{bookfigure}


\clearpage
\chapterformulaheading
\begin{formulasheetbox}
\begin{formulatable}
Position vector & $\mathbf r=x\mathbf e_x+y\mathbf e_y+z\mathbf e_z$\\
Magnitude & $\lVert\mathbf A\rVert=\sqrt{A_x^2+A_y^2+A_z^2}$\\
Dot product & $\mathbf A\cdot\mathbf B=A_xB_x+A_yB_y+A_zB_z
=\lVert\mathbf A\rVert\lVert\mathbf B\rVert\cos\theta$\\
Cross product magnitude &
$\lVert\mathbf A\times\mathbf B\rVert
=\lVert\mathbf A\rVert\lVert\mathbf B\rVert\sin\theta$\\
Unit vector & $\widehat{\mathbf A}=\mathbf A/\lVert\mathbf A\rVert$\\
\end{formulatable}
\end{formulasheetbox}

\chapternotationheading
\begin{notationbox}
\begin{notationtable}
$\mathbf r$ & Position vector & m\\
$\mathbf A,\mathbf B$ & General vectors & depends on quantity\\
$\mathbf e_x,\mathbf e_y,\mathbf e_z$ & Cartesian basis vectors & dimensionless\\
$\theta$ & Angle between vectors & rad\\
$\lVert\mathbf A\rVert$ & Magnitude of $\mathbf A$ & same as $\mathbf A$\\
\end{notationtable}
\end{notationbox}

\begin{warningbox}
A vector is not determined by magnitude alone. A speed of
$20\,\mathrm{m\,s^{-1}}$ is a scalar; a velocity of
$20\,\mathrm{m\,s^{-1}}$ east is a vector.
\end{warningbox}

\begin{chapterhistorical}[Cartesian geometry]
René Descartes and Pierre de Fermat developed analytic geometry in the
seventeenth century. Their central idea was to represent geometry using
coordinates and algebra, creating the language later used throughout mechanics,
electromagnetism, relativity, and quantum theory.
\end{chapterhistorical}

\chapteroutcomesheading
\begin{outcomesbox}
After studying this chapter, the reader should be able to:
\begin{itemize}
\item distinguish scalar and vector quantities;
\item resolve vectors into Cartesian components;
\item calculate magnitudes, dot products, and cross products;
\item use the right-hand rule;
\item explain why coordinate choices do not change physical laws.
\end{itemize}
\end{outcomesbox}

\chapterexercisesheading
\subsection*{Basic}
\begin{chapterexercise}[Vector magnitude]\difficulty{1}
For $\mathbf A=3\mathbf e_x-4\mathbf e_y+12\mathbf e_z$, calculate
$\lVert\mathbf A\rVert$ and the corresponding unit vector.
\end{chapterexercise}

\begin{chapterexercise}[Dot product]\difficulty{1}
For $\mathbf A=(2,-1,3)$ and $\mathbf B=(1,4,-2)$, calculate
$\mathbf A\cdot\mathbf B$ and determine whether the vectors are perpendicular.
\end{chapterexercise}

\subsection*{Intermediate}
\begin{chapterexercise}[Cross product and area]\difficulty{2}
Calculate $\mathbf A\times\mathbf B$ for
$\mathbf A=(1,2,0)$ and $\mathbf B=(0,1,3)$.
Interpret its magnitude geometrically.
\end{chapterexercise}

\subsection*{Advanced}
\begin{chapterexercise}[Coordinate independence]\difficulty{3}
Explain why the components of a vector change under a rotation of axes while
its magnitude does not.
\end{chapterexercise}

\chapterlookaheadheading
\begin{lookingaheadbox}
Vectors describe where an object is and in which direction it moves.
The next chapter makes these vectors functions of time and introduces
trajectory, velocity, and acceleration.
\[
\mathbf r
\longrightarrow
\mathbf r(t)
\longrightarrow
\mathbf v(t)
\longrightarrow
\mathbf a(t).
\]
\end{lookingaheadbox}


\chapter{Motion in Space and Time}
\chapteropeningart
{A trajectory is a curve through space}
{Position becomes a function of time. Its first and second derivatives produce
velocity and acceleration, the fundamental kinematic vectors.}
%%%%%%%%%%%%%%%%%%%%%%%%%%%%%%%%%%%%%%%%%%%%%%%%%%%%%%%%%%%%%%%%%%%%%%%%%%%%%%%%

\chapterquote
{``Nothing in physics is truly static. Every physical system evolves with time,
and understanding motion is the first step toward understanding nature.''}
{}

\begin{roadmapbox}
Position $\longrightarrow$ displacement $\longrightarrow$ velocity
$\longrightarrow$ acceleration $\longrightarrow$ trajectory.
\end{roadmapbox}

%%%%%%%%%%%%%%%%%%%%%%%%%%%%%%%%%%%%%%%%%%%%%%%%%%%%%%%%%%%%%%%%%%%%%%%%%%%%%%%%

\learningobjectives


\begin{itemize}
	\item understand the concept of motion;
	\item describe position as a function of time;
	\item distinguish between average and instantaneous quantities;
	\item calculate velocity and acceleration;
	\item understand derivative notation;
	\item describe motion using vectors;
	\item interpret graphs of position, velocity and acceleration;
	\item understand parametric curves and trajectories.
\end{itemize}

\newpage

%%%%%%%%%%%%%%%%%%%%%%%%%%%%%%%%%%%%%%%%%%%%%%%%%%%%%%%%%%%%%%%%%%%%%%%%%%%%%%%%

\section{What is Motion?}

One of the most fundamental questions in physics is:

\begin{center}
	\emph{What does it mean for an object to move?}
\end{center}

Suppose we observe an electron.

At time

\[
t_1
\]

its position is

\[
\mathbf r_1
=
(2,1,0)\ {\rm nm}.
\]

A short time later,

\[
t_2,
\]

its position becomes

\[
\mathbf r_2
=
(2.3,1.1,0)\ {\rm nm}.
\]

The position has changed.

Therefore,

\begin{keyidea}
Motion is the change of position with time.
\end{keyidea}

Motion is therefore not a property of space alone.

It requires two ingredients:

\begin{itemize}
	\item position,
	\item time.
\end{itemize}

%%%%%%%%%%%%%%%%%%%%%%%%%%%%%%%%%%%%%%%%%%%%%%%%%%%%%%%%%%%%%%%%%%%%%%%%%%%%%%%%

\section{Time}

Time is represented by

\[
t.
\]

Unlike spatial coordinates, time is an independent variable.

Instead of writing

\[
\mathbf r,
\]

we now write

\[
\boxed{
	\mathbf r(t).
}
\]

This notation means

\begin{quote}
	``The position depends on time.''
\end{quote}

%%%%%%%%%%%%%%%%%%%%%%%%%%%%%%%%%%%%%%%%%%%%%%%%%%%%%%%%%%%%%%%%%%%%%%%%%%%%%%%%

\section{Functions}

A function assigns exactly one output to every input.

Examples include

\[
x(t),
\]

\[
y(t),
\]

and

\[
z(t).
\]

Together they define the complete position vector

\[
\boxed{
	\mathbf r(t)
	=
	x(t)\hat{\mathbf x}
	+
	y(t)\hat{\mathbf y}
	+
	z(t)\hat{\mathbf z}.
}
\]

Every moving particle therefore traces a curve through space.

%%%%%%%%%%%%%%%%%%%%%%%%%%%%%%%%%%%%%%%%%%%%%%%%%%%%%%%%%%%%%%%%%%%%%%%%%%%%%%%%

\section{Trajectory}

The collection of all positions occupied by a moving particle is called its
trajectory.

Examples include

\begin{itemize}
	\item a straight line,
	\item a circle,
	\item an ellipse,
	\item a spiral,
	\item a completely irregular path.
\end{itemize}

Later, quantum mechanics will replace these classical trajectories by
wavefunctions, but the geometric picture remains an important intuition.

%%%%%%%%%%%%%%%%%%%%%%%%%%%%%%%%%%%%%%%%%%%%%%%%%%%%%%%%%%%%%%%%%%%%%%%%%%%%%%%%

\section{Average Velocity}

Suppose a particle moves

\[
\Delta x
=
10\ {\rm m}
\]

during

\[
\Delta t
=
2\ {\rm s}.
\]

Its average velocity is

\[
\boxed{
	v_{\rm avg}
	=
	\frac{\Delta x}{\Delta t}.
}
\]

For the above example,

\[
v_{\rm avg}
=
5\ {\rm m/s}.
\]

Average velocity tells us only what happened over an interval.

It does not describe the motion at every instant.

%%%%%%%%%%%%%%%%%%%%%%%%%%%%%%%%%%%%%%%%%%%%%%%%%%%%%%%%%%%%%%%%%%%%%%%%%%%%%%%%

\section{Instantaneous Velocity}

Imagine measuring the motion over smaller and smaller time intervals.

As

\[
\Delta t\rightarrow0,
\]

the average velocity becomes the instantaneous velocity.

Therefore,

\[
\boxed{
	v
	=
	\frac{dx}{dt}.
}
\]

In three dimensions,

\[
\boxed{
	\mathbf v
	=
	\frac{d\mathbf r}{dt}.
}
\]

Expanding into components,

\[
\boxed{
	\mathbf v
	=
	\frac{dx}{dt}\hat{\mathbf x}
	+
	\frac{dy}{dt}\hat{\mathbf y}
	+
	\frac{dz}{dt}\hat{\mathbf z}.
}
\]

Velocity therefore points in the direction of motion.

%%%%%%%%%%%%%%%%%%%%%%%%%%%%%%%%%%%%%%%%%%%%%%%%%%%%%%%%%%%%%%%%%%%%%%%%%%%%%%%%

\section{Speed}

Velocity is a vector.

Speed is a scalar.

The speed is simply the magnitude of the velocity vector,

\[
\boxed{
	v
	=
	|\mathbf v|.
}
\]

Thus,

\[
v
=
\sqrt{
	v_x^2
	+
	v_y^2
	+
	v_z^2
}.
\]

%%%%%%%%%%%%%%%%%%%%%%%%%%%%%%%%%%%%%%%%%%%%%%%%%%%%%%%%%%%%%%%%%%%%%%%%%%%%%%%%

\section{Acceleration}

Velocity itself may change.

The rate of change of velocity is called acceleration.

Average acceleration is

\[
\boxed{
	a_{\rm avg}
	=
	\frac{\Delta v}{\Delta t}.
}
\]

Instantaneous acceleration is

\[
\boxed{
	\mathbf a
	=
	\frac{d\mathbf v}{dt}.
}
\]

Since

\[
\mathbf v
=
\frac{d\mathbf r}{dt},
\]

we obtain

\[
\boxed{
	\mathbf a
	=
	\frac{d^2\mathbf r}{dt^2}.
}
\]

Acceleration measures how rapidly the velocity changes.

%%%%%%%%%%%%%%%%%%%%%%%%%%%%%%%%%%%%%%%%%%%%%%%%%%%%%%%%%%%%%%%%%%%%%%%%%%%%%%%%

\section{Dot Notation}

In analytical mechanics it is common to use Newton's dot notation.

Instead of writing

\[
\frac{dx}{dt},
\]

we write

\[
\boxed{
	\dot x.
}
\]

Similarly,

\[
\boxed{
	\ddot x
	=
	\frac{d^2x}{dt^2}.
}
\]

For vectors,

\[
\boxed{
	\dot{\mathbf r}
	=
	\mathbf v,
}
\]

and

\[
\boxed{
	\ddot{\mathbf r}
	=
	\mathbf a.
}
\]

This notation will be used extensively in the Lagrangian and Hamiltonian
formulations developed later in this book.

%%%%%%%%%%%%%%%%%%%%%%%%%%%%%%%%%%%%%%%%%%%%%%%%%%%%%%%%%%%%%%%%%%%%%%%%%%%%%%%%

\section{Uniform Circular Motion}

One of the most important examples of motion is uniform circular motion.

A particle moving around a circle of radius

\[
R
\]

with angular velocity

\[
\omega
\]

has position

\[
x(t)=R\cos(\omega t),
\]

\[
y(t)=R\sin(\omega t).
\]

Differentiating,

\[
v_x
=
-R\omega\sin(\omega t),
\]

\[
v_y
=
R\omega\cos(\omega t).
\]

The speed is constant,

\[
v
=
R\omega,
\]

while the direction changes continuously.

The acceleration always points toward the centre of the circle,

\[
\boxed{
	a
	=
	\frac{v^2}{R}.
}
\]

Later, we shall discover that electrons in a magnetic field naturally execute
cyclotron motion, which is closely related to this simple example.

%%%%%%%%%%%%%%%%%%%%%%%%%%%%%%%%%%%%%%%%%%%%%%%%%%%%%%%%%%%%%%%%%%%%%%%%%%%%%%%%


\begin{bookfigure}{A parametric trajectory with instantaneous velocity tangent
to the curve and acceleration changing the velocity vector.}
\begin{tikzpicture}[scale=1.0]
  \draw[physics axis] (-0.3,0) -- (9.2,0) node[right] {$x$};
  \draw[physics axis] (0,-0.3) -- (0,4.7) node[above] {$y$};
  \draw[trajectory]
    plot[smooth] coordinates {(0.4,0.5) (1.5,1.25) (2.6,2.45)
      (3.9,2.9) (5.1,2.3) (6.3,1.4) (7.5,1.65) (8.5,3.0)};
  \fill[visualgold] (5.1,2.3) circle (2.2pt);
  \node[below,visual label] at (5.1,2.3) {$\mathbf r(t)$};
  \draw[physics vector] (5.1,2.3) -- (6.65,1.38)
    node[right,visual label] {$\mathbf v(t)$};
  \draw[force vector] (5.1,2.3) -- (4.55,3.65)
    node[above,visual label] {$\mathbf a(t)$};
\end{tikzpicture}
\end{bookfigure}

\begin{visualguidebox}
The position vector identifies the point on the curve. Velocity is tangent to
the trajectory. Acceleration need not point along the direction of motion.
\end{visualguidebox}


\section{Motion as Parametric Curves}

The position vector

\[
\mathbf r(t)
\]

defines a parametric curve.

The velocity vector is tangent to this curve,

\[
\mathbf v
=
\frac{d\mathbf r}{dt},
\]

while the acceleration describes how the tangent vector changes.

These geometric ideas later become central in differential geometry and
general relativity.

%%%%%%%%%%%%%%%%%%%%%%%%%%%%%%%%%%%%%%%%%%%%%%%%%%%%%%%%%%%%%%%%%%%%%%%%%%%%%%%%

\section{Connection to Future Chapters}

The concepts introduced here appear everywhere in physics.

Newton's Second Law uses

\[
\mathbf a.
\]

The Lagrangian depends upon

\[
\mathbf r
\quad\text{and}\quad
\dot{\mathbf r}.
\]

The Hamiltonian depends upon

\[
\mathbf r
\quad\text{and}\quad
\mathbf p.
\]

Quantum mechanics introduces

\[
\psi(\mathbf r,t),
\]

which also depends on position and time.

Thus, understanding motion is an essential prerequisite for all later chapters.

%%%%%%%%%%%%%%%%%%%%%%%%%%%%%%%%%%%%%%%%%%%%%%%%%%%%%%%%%%%%%%%%%%%%%%%%%%%%%%%%

\section{Summary}

In this chapter we introduced

\begin{itemize}
	\item motion,
	\item time,
	\item functions,
	\item trajectories,
	\item average velocity,
	\item instantaneous velocity,
	\item speed,
	\item acceleration,
	\item dot notation,
	\item circular motion,
	\item parametric curves.
\end{itemize}

The next chapter develops the mathematical tools needed to describe changing
quantities rigorously. We shall introduce limits, derivatives, partial
derivatives and the differential operators that form the language of
electromagnetism and quantum mechanics.


%%%%%%%%%%%%%%%%%%%%%%%%%%%%%%%%%%%%%%%%%%%%%%%%%%%%%%%%%%%%%%%%%%%%%%%%%%%%%%%%

\clearpage
\chapterformulaheading
\begin{formulasheetbox}
\begin{formulatable}
Displacement & $\Delta\mathbf r=\mathbf r(t_2)-\mathbf r(t_1)$\\
Average velocity & $\overline{\mathbf v}=\Delta\mathbf r/\Delta t$\\
Instantaneous velocity & $\mathbf v=d\mathbf r/dt$\\
Speed & $v=\lVert\mathbf v\rVert$\\
Acceleration & $\mathbf a=d\mathbf v/dt=d^2\mathbf r/dt^2$\\
Uniform circular motion & $\lVert\mathbf a\rVert=v^2/R=\omega^2R$\\
\end{formulatable}
\end{formulasheetbox}

\chapternotationheading
\begin{notationbox}
\begin{notationtable}
$t$ & Time & s\\
$\mathbf r(t)$ & Position as a function of time & m\\
$\Delta\mathbf r$ & Displacement & m\\
$\mathbf v$ & Velocity & m\,s$^{-1}$\\
$v$ & Speed & m\,s$^{-1}$\\
$\mathbf a$ & Acceleration & m\,s$^{-2}$\\
$\omega$ & Angular velocity & rad\,s$^{-1}$\\
\end{notationtable}
\end{notationbox}

\begin{physical}
Velocity is tangent to the trajectory, whereas acceleration describes how that
tangent vector changes. Consequently, an object may accelerate even when its
speed remains constant, as in uniform circular motion.
\end{physical}

\begin{warningbox}
Average velocity and instantaneous velocity are different concepts.
The average uses a finite time interval; the instantaneous value is obtained
from the limiting derivative at one time.
\end{warningbox}

\begin{chapterhistorical}[Galileo and mathematical motion]
Galileo's studies of falling bodies and projectiles helped replace qualitative
Aristotelian descriptions with mathematical laws involving distance, time, and
acceleration. This transition prepared the way for Newtonian mechanics.
\end{chapterhistorical}

\chapteroutcomesheading
\begin{outcomesbox}
The reader should now be able to:
\begin{itemize}
\item interpret a trajectory $\mathbf r(t)$;
\item distinguish distance, displacement, speed, and velocity;
\item obtain velocity and acceleration by differentiation;
\item describe circular and parametric motion;
\item connect kinematics to later Lagrangian and Hamiltonian descriptions.
\end{itemize}
\end{outcomesbox}

\chapterexercisesheading
\subsection*{Basic}
\begin{chapterexercise}[One-dimensional motion]\difficulty{1}
Let $x(t)=4t^2-2t+1$. Determine the velocity and acceleration.
\end{chapterexercise}

\begin{chapterexercise}[Average and instantaneous velocity]\difficulty{2}
For $\mathbf r(t)=(t^2,2t,t^3)$, calculate the average velocity from
$t=1$ to $t=2$ and the instantaneous velocity at $t=2$.
\end{chapterexercise}

\subsection*{Intermediate}
\begin{chapterexercise}[Circular trajectory]\difficulty{2}
For $\mathbf r(t)=R(\cos\omega t,\sin\omega t,0)$, derive
$\mathbf v(t)$ and $\mathbf a(t)$ and show that
$\mathbf a=-\omega^2\mathbf r$.
\end{chapterexercise}

\subsection*{Programming}
\begin{chapterexercise}[Numerical trajectory]\difficulty{3}
Sample a three-dimensional parametric trajectory at discrete times, estimate
its velocity using finite differences, and compare the estimate with the
analytic derivative.
\end{chapterexercise}

\chapterlookaheadheading
\begin{lookingaheadbox}
Kinematics uses derivatives, but a rigorous understanding of derivatives
requires limits and multivariable calculus. The next chapter develops these
tools and introduces $\nabla$, divergence, curl, and the Laplacian.
\end{lookingaheadbox}


\chapter{Calculus and Differential Operators}
\chapteropeningart
{Local change reveals global structure}
{Derivatives measure change; gradient, divergence, curl, and the Laplacian
describe the local geometry of scalar and vector fields.}
%%%%%%%%%%%%%%%%%%%%%%%%%%%%%%%%%%%%%%%%%%%%%%%%%%%%%%%%%%%%%%%%%%%%%%%%%%%%%%%%

\chapterquote
{``Nature is written in the language of mathematics, and calculus is the language
used to describe continuous change.''}
{}

\begin{roadmapbox}
Functions $\longrightarrow$ limits $\longrightarrow$ derivatives
$\longrightarrow$ partial derivatives $\longrightarrow$
gradient, divergence, curl, and Laplacian.
\end{roadmapbox}

%%%%%%%%%%%%%%%%%%%%%%%%%%%%%%%%%%%%%%%%%%%%%%%%%%%%%%%%%%%%%%%%%%%%%%%%%%%%%%%%

\learningobjectives

\begin{itemize}
	\item understand why calculus was invented;
	\item explain the concept of a limit;
	\item define the derivative from first principles;
	\item distinguish average and instantaneous rates of change;
	\item calculate first and second derivatives;
	\item understand partial derivatives;
	\item distinguish total and partial derivatives;
	\item understand the differential operator $\nabla$;
	\item calculate gradients, divergences, curls and Laplacians;
	\item recognize these operators in physics.
\end{itemize}

\newpage

%%%%%%%%%%%%%%%%%%%%%%%%%%%%%%%%%%%%%%%%%%%%%%%%%%%%%%%%%%%%%%%%%%%%%%%%%%%%%%%%

\section{Why Calculus?}

Imagine observing an electron moving along a straight line.

At different moments we measure

\[
\begin{array}{c|c}
	t\;(\mathrm{s}) & x\;(\mathrm{m})\\
	\hline
	0 & 0\\
	1 & 2\\
	2 & 5\\
	3 & 9
\end{array}
\]

The particle is clearly moving.

More importantly, its speed is changing.

The natural question is

\begin{center}
	\emph{How can we describe continuous change mathematically?}
\end{center}

The answer is calculus.

Calculus was independently developed by Isaac Newton and Gottfried Wilhelm
Leibniz during the seventeenth century to describe continuously changing
physical systems.

%%%%%%%%%%%%%%%%%%%%%%%%%%%%%%%%%%%%%%%%%%%%%%%%%%%%%%%%%%%%%%%%%%%%%%%%%%%%%%%%

\section{Limits}

Before defining a derivative, we must understand the idea of a limit.

Consider

\[
f(x)=\frac{x^2-1}{x-1}.
\]

At

\[
x=1,
\]

the denominator becomes zero.

At first sight the function appears undefined.

However,

\[
x^2-1=(x-1)(x+1),
\]

therefore

\[
f(x)=x+1,
\qquad x\neq1.
\]

As

\[
x\rightarrow1,
\]

the function approaches

\[
2.
\]

Therefore

\[
\boxed{
	\lim_{x\rightarrow1}
	\frac{x^2-1}{x-1}
	=2.
}
\]

A limit describes the behaviour of a function as the input approaches
a particular value.

%%%%%%%%%%%%%%%%%%%%%%%%%%%%%%%%%%%%%%%%%%%%%%%%%%%%%%%%%%%%%%%%%%%%%%%%%%%%%%%%

\section{Average Rate of Change}

Suppose the value of a function changes from

\[
f(x)
\]

to

\[
f(x+\Delta x).
\]

The average rate of change is

\[
\boxed{
	\frac{\Delta f}{\Delta x}
	=
	\frac{f(x+\Delta x)-f(x)}
	{\Delta x}.
}
\]

This is exactly the average velocity introduced in the previous chapter.

%%%%%%%%%%%%%%%%%%%%%%%%%%%%%%%%%%%%%%%%%%%%%%%%%%%%%%%%%%%%%%%%%%%%%%%%%%%%%%%%

\section{The Derivative}

To obtain the instantaneous rate of change, we make the interval
smaller and smaller.

Mathematically,

\[
\Delta x\rightarrow0.
\]

The derivative is therefore defined as

\[
\boxed{
	\frac{df}{dx}
	=
	\lim_{\Delta x\rightarrow0}
	\frac{f(x+\Delta x)-f(x)}
	{\Delta x}.
}
\]

\begin{keyidea}
The derivative is an instantaneous rate of change. Geometrically, it is the
slope of the tangent line to the graph.
\end{keyidea}

%%%%%%%%%%%%%%%%%%%%%%%%%%%%%%%%%%%%%%%%%%%%%%%%%%%%%%%%%%%%%%%%%%%%%%%%%%%%%%%%

\section{Physical Meaning of the Derivative}

If

\[
x=x(t),
\]

then

\[
\boxed{
	v=\frac{dx}{dt}
}
\]

is the velocity.

Differentiating once again,

\[
\boxed{
	a=\frac{d^2x}{dt^2},
}
\]

which represents the acceleration.

Thus,

\[
\boxed{
	\text{Derivative}
	=
	\text{Rate of change.}
}
\]

%%%%%%%%%%%%%%%%%%%%%%%%%%%%%%%%%%%%%%%%%%%%%%%%%%%%%%%%%%%%%%%%%%%%%%%%%%%%%%%%

\section{Higher Derivatives}

The first derivative measures how rapidly a quantity changes.

The second derivative measures how rapidly the first derivative changes.

Examples include

\[
\frac{d^2x}{dt^2},
\]

\[
\frac{d^3x}{dt^3},
\]

and higher-order derivatives.

These appear naturally in mechanics, elasticity and wave propagation.

%%%%%%%%%%%%%%%%%%%%%%%%%%%%%%%%%%%%%%%%%%%%%%%%%%%%%%%%%%%%%%%%%%%%%%%%%%%%%%%%

\section{Partial Derivatives}

Many physical quantities depend on several variables.

Suppose

\[
f(x,y,z)=x^2+y^2+z.
\]

Differentiating only with respect to

\[
x,
\]

while treating

\[
y,z
\]

as constants gives

\[
\boxed{
	\frac{\partial f}{\partial x}=2x.
}
\]

Similarly,

\[
\frac{\partial f}{\partial y}=2y,
\]

and

\[
\frac{\partial f}{\partial z}=1.
\]

Partial derivatives describe how a function changes along one coordinate
direction while all other coordinates remain fixed.

%%%%%%%%%%%%%%%%%%%%%%%%%%%%%%%%%%%%%%%%%%%%%%%%%%%%%%%%%%%%%%%%%%%%%%%%%%%%%%%%

\section{Total Derivative}

Suppose

\[
f=f(x(t),y(t)).
\]

Now every variable depends upon time.

The total derivative becomes

\[
\boxed{
	\frac{df}{dt}
	=
	\frac{\partial f}{\partial x}
	\frac{dx}{dt}
	+
	\frac{\partial f}{\partial y}
	\frac{dy}{dt}.
}
\]

This is the multivariable chain rule.

%%%%%%%%%%%%%%%%%%%%%%%%%%%%%%%%%%%%%%%%%%%%%%%%%%%%%%%%%%%%%%%%%%%%%%%%%%%%%%%%

\section{The Differential Operator}

The symbol

\[
\boxed{
	\nabla
}
\]

is called

\begin{itemize}
	\item the del operator,
	\item the nabla operator,
	\item the vector differential operator.
\end{itemize}

It is defined as

\[
\boxed{
	\nabla
	=
	\hat{\mathbf x}\frac{\partial}{\partial x}
	+
	\hat{\mathbf y}\frac{\partial}{\partial y}
	+
	\hat{\mathbf z}\frac{\partial}{\partial z}.
}
\]

Unlike an ordinary vector,

\[
\nabla
\]

contains differentiation operators instead of numerical components.

%%%%%%%%%%%%%%%%%%%%%%%%%%%%%%%%%%%%%%%%%%%%%%%%%%%%%%%%%%%%%%%%%%%%%%%%%%%%%%%%

\section{Gradient}


\begin{bookfigure}{Gradient vectors are perpendicular to level curves and point
toward the steepest local increase of the scalar field.}
\begin{tikzpicture}[scale=1.0]
  \foreach \r/\lab in {0.8/1,1.5/2,2.25/3,3.05/4}{
    \draw[visualcyan!75,line width=0.9pt] (0,0) ellipse ({1.35*\r} and {\r});
    \node[visual label] at ({1.35*\r},0.2) {$f=\lab$};
  }
  \foreach \a in {20,65,115,165,215,265,315}{
    \draw[field vector] ({1.3*cos(\a)},{0.95*sin(\a)})
      -- ({2.25*cos(\a)},{1.65*sin(\a)});
  }
  \fill[visualgold] (1.15,0.52) circle (2pt);
  \draw[force vector] (1.15,0.52) -- (2.4,1.08)
    node[above right,visual label] {$\nabla f$};
  \node[font=\small,align=center,text width=4.2cm] at (7.2,0.4)
    {Level curves connect points of equal field value.\\[2mm]
     The gradient crosses them orthogonally.};
\end{tikzpicture}
\end{bookfigure}


If

\[
f(x,y,z)
\]

is a scalar field,

then

\[
\boxed{
	\nabla f
	=
	\left(
	\frac{\partial f}{\partial x},
	\frac{\partial f}{\partial y},
	\frac{\partial f}{\partial z}
	\right).
}
\]

The gradient points toward the direction of greatest increase of the function.

For example,

\[
T(x,y)
\]

may represent temperature.

Then

\[
\nabla T
\]

points in the direction where the temperature rises most rapidly.

%%%%%%%%%%%%%%%%%%%%%%%%%%%%%%%%%%%%%%%%%%%%%%%%%%%%%%%%%%%%%%%%%%%%%%%%%%%%%%%%

\section{Divergence}

For a vector field

\[
\mathbf A
=
(A_x,A_y,A_z),
\]

the divergence is

\[
\boxed{
	\nabla\cdot\mathbf A
	=
	\frac{\partial A_x}{\partial x}
	+
	\frac{\partial A_y}{\partial y}
	+
	\frac{\partial A_z}{\partial z}.
}
\]

The divergence measures the strength of local sources or sinks.

%%%%%%%%%%%%%%%%%%%%%%%%%%%%%%%%%%%%%%%%%%%%%%%%%%%%%%%%%%%%%%%%%%%%%%%%%%%%%%%%

\section{Curl}

The curl is

\[
\boxed{
	\nabla\times\mathbf A.
}
\]

Explicitly,

\[
\nabla\times\mathbf A
=
\begin{pmatrix}
	\partial_yA_z-\partial_zA_y\\
	\partial_zA_x-\partial_xA_z\\
	\partial_xA_y-\partial_yA_x
\end{pmatrix}.
\]

The curl measures local rotation.

In electromagnetism,

\[
\boxed{
	\mathbf B
	=
	\nabla\times\mathbf A.
}
\]

%%%%%%%%%%%%%%%%%%%%%%%%%%%%%%%%%%%%%%%%%%%%%%%%%%%%%%%%%%%%%%%%%%%%%%%%%%%%%%%%

\section{Laplacian}

Applying the divergence to the gradient gives

\[
\boxed{
	\nabla^2f
	=
	\nabla\cdot(\nabla f).
}
\]

In Cartesian coordinates,

\[
\boxed{
	\nabla^2f
	=
	\frac{\partial^2f}{\partial x^2}
	+
	\frac{\partial^2f}{\partial y^2}
	+
	\frac{\partial^2f}{\partial z^2}.
}
\]

The Laplacian measures the local curvature of a scalar field.

%%%%%%%%%%%%%%%%%%%%%%%%%%%%%%%%%%%%%%%%%%%%%%%%%%%%%%%%%%%%%%%%%%%%%%%%%%%%%%%%

\section{Important Vector Identities}

Several identities appear repeatedly in physics.

\[
\boxed{
	\nabla\times(\nabla f)=0.
}
\]

The curl of a gradient always vanishes.

Similarly,

\[
\boxed{
	\nabla\cdot(\nabla\times\mathbf A)=0.
}
\]

The divergence of a curl is always zero.

These identities play a central role in Maxwell's equations.

%%%%%%%%%%%%%%%%%%%%%%%%%%%%%%%%%%%%%%%%%%%%%%%%%%%%%%%%%%%%%%%%%%%%%%%%%%%%%%%%

\section{Applications in Physics}

Calculus appears throughout physics.

Examples include

\[
v=\frac{dx}{dt},
\]

\[
a=\frac{d^2x}{dt^2},
\]

\[
\mathbf E=-\nabla\phi,
\]

\[
\mathbf B=\nabla\times\mathbf A,
\]

\[
\hat{\mathbf p}
=
-i\hbar\nabla,
\]

and

\[
\hat H
=
-\frac{\hbar^2}{2m}\nabla^2+V.
\]

Thus a single mathematical idea—the derivative—connects mechanics,
electromagnetism and quantum mechanics.

%%%%%%%%%%%%%%%%%%%%%%%%%%%%%%%%%%%%%%%%%%%%%%%%%%%%%%%%%%%%%%%%%%%%%%%%%%%%%%%%

\section{Summary}

In this chapter we introduced

\begin{itemize}
	\item limits;
	\item derivatives;
	\item higher derivatives;
	\item partial derivatives;
	\item total derivatives;
	\item the differential operator $\nabla$;
	\item the gradient;
	\item divergence;
	\item curl;
	\item the Laplacian.
\end{itemize}

These operators form the mathematical foundation of nearly every field of
theoretical physics. In the next chapter we shall apply them to classical
mechanics and derive Newton's laws of motion in a mathematically rigorous way.


%%%%%%%%%%%%%%%%%%%%%%%%%%%%%%%%%%%%%%%%%%%%%%%%%%%%%%%%%%%%%%%%%%%%%%%%%%%%%%%%

\clearpage
\chapterformulaheading
\begin{formulasheetbox}
\begin{formulatable}
Derivative & $f'(x)=\displaystyle\lim_{h\to0}\frac{f(x+h)-f(x)}{h}$\\
Gradient & $\nabla f=
\mathbf e_x\partial_x f+\mathbf e_y\partial_y f+\mathbf e_z\partial_z f$\\
Divergence & $\nabla\cdot\mathbf A=
\partial_xA_x+\partial_yA_y+\partial_zA_z$\\
Curl & $\nabla\times\mathbf A$\\
Laplacian & $\nabla^2 f=
\partial_x^2f+\partial_y^2f+\partial_z^2f$\\
Total derivative & $\displaystyle\frac{df}{dt}
=\partial_t f+\dot x\,\partial_xf+\dot y\,\partial_yf+\dot z\,\partial_zf$\\
\end{formulatable}
\end{formulasheetbox}

\chapternotationheading
\begin{notationbox}
\begin{notationtable}
$d/dx$ & Ordinary derivative operator & inverse unit of $x$\\
$\partial/\partial x$ & Partial derivative operator & inverse unit of $x$\\
$\nabla$ & Vector differential operator & m$^{-1}$ in space\\
$\nabla f$ & Gradient of scalar field & unit of $f$ per m\\
$\nabla\cdot\mathbf A$ & Divergence of vector field & unit of $\mathbf A$ per m\\
$\nabla\times\mathbf A$ & Curl of vector field & unit of $\mathbf A$ per m\\
$\nabla^2$ & Laplacian & m$^{-2}$ in space\\
\end{notationtable}
\end{notationbox}

\begin{physical}
The gradient points toward the steepest local increase of a scalar field.
Divergence measures local source or sink behaviour. Curl measures local
circulation. The Laplacian compares a value with its spatial neighbourhood.
These interpretations reappear in electrostatics, fluid mechanics, diffusion,
and the Schrödinger equation.
\end{physical}

\begin{warningbox}
A partial derivative changes one independent variable while holding the others
fixed. A total derivative follows the complete dependence of a quantity,
including indirect dependence through variables such as $x(t)$, $y(t)$,
and $z(t)$.
\end{warningbox}

\begin{chapterhistorical}[Newton and Leibniz]
Isaac Newton and Gottfried Wilhelm Leibniz developed calculus independently in
the seventeenth century. Newton emphasized changing physical quantities, while
Leibniz introduced notation that remains central to modern analysis.
\end{chapterhistorical}

\chapteroutcomesheading
\begin{outcomesbox}
The reader should now be able to:
\begin{itemize}
\item define a derivative using a limit;
\item distinguish ordinary, partial, and total derivatives;
\item compute gradients, divergence, curl, and Laplacians;
\item give geometric and physical interpretations of these operators;
\item recognize the same operators in classical and quantum field equations.
\end{itemize}
\end{outcomesbox}

\chapterexercisesheading
\subsection*{Basic}
\begin{chapterexercise}[Derivative from first principles]\difficulty{2}
Use the limit definition to derive $d(x^2)/dx=2x$.
\end{chapterexercise}

\begin{chapterexercise}[Gradient]\difficulty{1}
Calculate $\nabla f$ for $f(x,y,z)=x^2y+yz^2$.
\end{chapterexercise}

\subsection*{Intermediate}
\begin{chapterexercise}[Divergence and curl]\difficulty{2}
For $\mathbf A=(xy,yz,zx)$, calculate
$\nabla\cdot\mathbf A$ and $\nabla\times\mathbf A$.
\end{chapterexercise}

\subsection*{Advanced}
\begin{chapterexercise}[Laplacian and waves]\difficulty{3}
Calculate $\nabla^2 e^{i\mathbf k\cdot\mathbf r}$ and interpret the result.
\end{chapterexercise}

\chapterlookaheadheading
\begin{lookingaheadbox}
The next chapter applies derivatives to momentum and acceleration. Kinematics
becomes dynamics when forces determine how motion changes:
\[
\mathbf a
\longrightarrow
m\mathbf a
\longrightarrow
\mathbf F.
\]
\end{lookingaheadbox}


\chapter{Newton's Laws of Motion}
\chapteropeningart
{Forces determine the evolution of motion}
{Newton's laws connect interactions to momentum and acceleration, turning
kinematics into predictive dynamics.}
%%%%%%%%%%%%%%%%%%%%%%%%%%%%%%%%%%%%%%%%%%%%%%%%%%%%%%%%%%%%%%%%%%%%%%%%%%%%%%%%

\chapterquote
{``To every action there is always opposed an equal reaction.''}
{Isaac Newton, \emph{Philosophiae Naturalis Principia Mathematica} (1687)}

\begin{roadmapbox}
Inertia $\longrightarrow$ force $\longrightarrow$ momentum
$\longrightarrow$ equations of motion $\longrightarrow$ free-body analysis.
\end{roadmapbox}

%%%%%%%%%%%%%%%%%%%%%%%%%%%%%%%%%%%%%%%%%%%%%%%%%%%%%%%%%%%%%%%%%%%%%%%%%%%%%%%%

\learningobjectives

\begin{itemize}
	\item understand the concept of force;
	\item distinguish between mass and weight;
	\item explain Newton's three laws of motion;
	\item calculate the motion of particles under applied forces;
	\item understand inertia and momentum;
	\item appreciate why Newtonian mechanics is an approximation to relativity
	and quantum mechanics.
\end{itemize}

%%%%%%%%%%%%%%%%%%%%%%%%%%%%%%%%%%%%%%%%%%%%%%%%%%%%%%%%%%%%%%%%%%%%%%%%%%%%%%%%

\section{Classical Mechanics}

Classical mechanics describes the motion of macroscopic objects.

Examples include

\begin{itemize}
	\item planets,
	\item automobiles,
	\item projectiles,
	\item pendulums,
	\item satellites.
\end{itemize}

For objects moving much slower than the speed of light and much larger than
atoms, Newton's laws provide an extremely accurate description of nature.

%%%%%%%%%%%%%%%%%%%%%%%%%%%%%%%%%%%%%%%%%%%%%%%%%%%%%%%%%%%%%%%%%%%%%%%%%%%%%%%%

\section{The Concept of Force}

In everyday language we often speak of pushing or pulling an object.

Physics introduces the more general concept of a force.

A force is any interaction capable of changing the motion of an object.

Examples include

\begin{itemize}
	\item gravity,
	\item electric forces,
	\item magnetic forces,
	\item friction,
	\item tension,
	\item spring forces.
\end{itemize}

The SI unit of force is the newton (N),

\[
1\;\mathrm N
=
1\;\mathrm{kg\,m/s^2}.
\]

%%%%%%%%%%%%%%%%%%%%%%%%%%%%%%%%%%%%%%%%%%%%%%%%%%%%%%%%%%%%%%%%%%%%%%%%%%%%%%%%

\section{Newton's First Law}

Newton's first law states

\begin{quote}
	Every body remains at rest or continues moving with constant velocity unless
	acted upon by a net external force.
\end{quote}

This law introduces the concept of inertia.

Objects naturally resist changes in their state of motion.

%%%%%%%%%%%%%%%%%%%%%%%%%%%%%%%%%%%%%%%%%%%%%%%%%%%%%%%%%%%%%%%%%%%%%%%%%%%%%%%%

\section{Inertial Reference Frames}

Newton's laws are valid only in inertial reference frames.

An inertial frame is one in which an isolated object moves with constant velocity.

Accelerating reference frames require additional fictitious forces such as

\begin{itemize}
	\item centrifugal force,
	\item Coriolis force.
\end{itemize}

%%%%%%%%%%%%%%%%%%%%%%%%%%%%%%%%%%%%%%%%%%%%%%%%%%%%%%%%%%%%%%%%%%%%%%%%%%%%%%%%

\section{Mass}

Mass measures inertia.

Objects with larger mass require larger forces to produce the same acceleration.

Mass should not be confused with weight. Mass measures inertia, whereas
weight is the gravitational force acting on that mass.

\begin{mistakebox}
Mass is measured in kilograms; weight is a force measured in newtons.
\end{mistakebox}

%%%%%%%%%%%%%%%%%%%%%%%%%%%%%%%%%%%%%%%%%%%%%%%%%%%%%%%%%%%%%%%%%%%%%%%%%%%%%%%%

\section{Newton's Second Law}


\begin{bookfigure}{Free-body diagram for a block on a horizontal surface. Only
forces acting on the selected body are shown.}
\begin{tikzpicture}[x=1cm,y=1cm]
  \draw[line width=1pt,draw=visualink] (-0.5,0) -- (8.5,0);
  \foreach \x in {-0.4,0.1,...,8.2}
    \draw[visualgrid] (\x,0) -- ++(-0.25,-0.25);
  \fill[softblue] (2.7,0) rectangle (5.6,1.75);
  \draw[visualblue,line width=1pt] (2.7,0) rectangle (5.6,1.75);
  \node at (4.15,0.88) {$m$};

  \draw[force vector] (4.15,1.75) -- (4.15,3.45)
    node[above,visual label] {$\mathbf N$};
  \draw[force vector] (4.15,0) -- (4.15,-1.65)
    node[below,visual label] {$m\mathbf g$};
  \draw[force vector] (5.6,0.9) -- (7.55,0.9)
    node[right,visual label] {$\mathbf F_{\mathrm{applied}}$};
  \draw[force vector] (2.7,0.55) -- (1.25,0.55)
    node[left,visual label] {$\mathbf f$};

  \node[concept node,text width=3.2cm] at (9.8,2.6)
    {$\sum\mathbf F=m\mathbf a$\\Vector sum first; acceleration second.};
\end{tikzpicture}
\end{bookfigure}

\begin{visualguidebox}
The normal force and weight cancel vertically when there is no vertical
acceleration. The horizontal net force determines the horizontal acceleration.
\end{visualguidebox}


The most famous equation in classical mechanics is

\[
\boxed{
	\mathbf F
	=
	m\mathbf a.
}
\]

Since

\[
\mathbf a
=
\frac{d^2\mathbf r}{dt^2},
\]

we may also write

\[
\boxed{
	m
	\frac{d^2\mathbf r}{dt^2}
	=
	\mathbf F.
}
\]

This is a second-order differential equation.

Given the force, solving this equation determines the particle's motion.

%%%%%%%%%%%%%%%%%%%%%%%%%%%%%%%%%%%%%%%%%%%%%%%%%%%%%%%%%%%%%%%%%%%%%%%%%%%%%%%%

\section{Momentum}

Momentum is defined by

\[
\boxed{
	\mathbf p
	=
	m\mathbf v.
}
\]

Using momentum,

Newton's second law becomes

\[
\boxed{
	\mathbf F
	=
	\frac{d\mathbf p}{dt}.
}
\]

This form remains valid even when the mass changes.

%%%%%%%%%%%%%%%%%%%%%%%%%%%%%%%%%%%%%%%%%%%%%%%%%%%%%%%%%%%%%%%%%%%%%%%%%%%%%%%%

\section{Newton's Third Law}

Newton's third law states

\begin{quote}
	For every action there exists an equal and opposite reaction.
\end{quote}

If object A exerts a force upon object B,

then object B simultaneously exerts an equal and opposite force upon object A.

Mathematically,

\[
\boxed{
	\mathbf F_{AB}
	=
	-
	\mathbf F_{BA}.
}
\]

%%%%%%%%%%%%%%%%%%%%%%%%%%%%%%%%%%%%%%%%%%%%%%%%%%%%%%%%%%%%%%%%%%%%%%%%%%%%%%%%

\section{Free-Body Diagrams}

To analyse motion, all forces acting on an object should first be identified.

Typical forces include

\begin{itemize}
	\item gravity,
	\item normal force,
	\item friction,
	\item tension,
	\item external applied forces.
\end{itemize}

The vector sum of these forces determines the acceleration.

%%%%%%%%%%%%%%%%%%%%%%%%%%%%%%%%%%%%%%%%%%%%%%%%%%%%%%%%%%%%%%%%%%%%%%%%%%%%%%%%

\section{Worked Example}

\begin{examplebox}
A force \(\mathbf F=(10,0,0)\,\mathrm N\) acts on a \(2\,\mathrm{kg}\) mass.
Find the acceleration.
\end{examplebox}

\begin{solutionbox}
Newton's second law gives
\[
\mathbf a=\frac{\mathbf F}{m}
=\frac{(10,0,0)}{2}
=(5,0,0)\,\mathrm{m/s^2}.
\]
\end{solutionbox}

%%%%%%%%%%%%%%%%%%%%%%%%%%%%%%%%%%%%%%%%%%%%%%%%%%%%%%%%%%%%%%%%%%%%%%%%%%%%%%%%

\section{Limitations of Newtonian Mechanics}

Although remarkably successful, Newton's theory is not universal.

For

\begin{itemize}
	\item velocities approaching the speed of light,
	\item atomic dimensions,
	\item strong gravitational fields,
\end{itemize}

more advanced theories become necessary.

These include

\begin{itemize}
	\item Special Relativity,
	\item General Relativity,
	\item Quantum Mechanics.
\end{itemize}

Nevertheless, Newtonian mechanics remains the starting point of nearly every
modern physical theory.

%%%%%%%%%%%%%%%%%%%%%%%%%%%%%%%%%%%%%%%%%%%%%%%%%%%%%%%%%%%%%%%%%%%%%%%%%%%%%%%%

\section{Connection to the Next Chapter}

Newton's equation

\[
m\mathbf a=\mathbf F
\]

describes motion once the force is known.

In the next chapter we shall study the most important forces in nature,
beginning with gravitational, electric and magnetic interactions.

Eventually we shall derive the Lorentz force,

\begin{keyequation}{Lorentz force}
\[
\mathbf F=q(\mathbf E+\mathbf v\times\mathbf B).
\]
\end{keyequation}

which forms the classical foundation of the Quantum Hall Effect.

%%%%%%%%%%%%%%%%%%%%%%%%%%%%%%%%%%%%%%%%%%%%%%%%%%%%%%%%%%%%%%%%%%%%%%%%%%%%%%%%

\section{Summary}

This chapter introduced

\begin{itemize}
	\item force;
	\item inertia;
	\item inertial frames;
	\item mass;
	\item Newton's three laws;
	\item momentum;
	\item free-body diagrams;
	\item limitations of classical mechanics.
\end{itemize}

These ideas provide the bridge between elementary kinematics and the
electromagnetic forces that govern charged particles.


\clearpage
\chapterformulaheading
\begin{formulasheetbox}
\begin{formulatable}
Momentum & $\mathbf p=m\mathbf v$\\
Newton's second law & $\mathbf F_{\mathrm{net}}=d\mathbf p/dt$\\
Constant-mass form & $\mathbf F_{\mathrm{net}}=m\mathbf a$\\
Weight near Earth's surface & $\mathbf W=m\mathbf g$\\
Action--reaction pair & $\mathbf F_{A\to B}=-\mathbf F_{B\to A}$\\
Impulse & $\mathbf J=\displaystyle\int_{t_1}^{t_2}\mathbf F\,dt=\Delta\mathbf p$\\
\end{formulatable}
\end{formulasheetbox}

\chapternotationheading
\begin{notationbox}
\begin{notationtable}
$m$ & Inertial mass & kg\\
$\mathbf F$ & Force & N\\
$\mathbf F_{\mathrm{net}}$ & Vector sum of forces & N\\
$\mathbf p$ & Momentum & kg\,m\,s$^{-1}$\\
$\mathbf g$ & Gravitational acceleration & m\,s$^{-2}$\\
$\mathbf J$ & Impulse & N\,s\\
\end{notationtable}
\end{notationbox}

\begin{physical}
Newton's second law does not merely state that force causes motion. It states
that net force controls the rate of change of momentum. The familiar equation
$\mathbf F=m\mathbf a$ is the constant-mass special case.
\end{physical}

\begin{warningbox}
Mass and weight are not interchangeable. Mass measures inertia and is expressed
in kilograms. Weight is a gravitational force and is expressed in newtons.
Also, the two forces in an action--reaction pair act on different bodies, so
they do not cancel in a free-body diagram of one body.
\end{warningbox}

\begin{chapterhistorical}[The \emph{Principia}]
Newton published the \emph{Philosophiae Naturalis Principia Mathematica} in
1687. It unified terrestrial and celestial motion using a small set of
mathematical laws and became the foundation of classical mechanics.
\end{chapterhistorical}

\chapteroutcomesheading
\begin{outcomesbox}
The reader should now be able to:
\begin{itemize}
\item identify inertial frames;
\item distinguish mass, weight, momentum, and force;
\item apply Newton's three laws;
\item construct and interpret free-body diagrams;
\item derive equations of motion from a specified net force;
\item state the limits of Newtonian mechanics.
\end{itemize}
\end{outcomesbox}

\chapterexercisesheading
\subsection*{Basic}
\begin{chapterexercise}[Net force]\difficulty{1}
A $4\,\mathrm{kg}$ body experiences forces
$12\,\mathrm N$ east and $5\,\mathrm N$ west. Determine its acceleration.
\end{chapterexercise}

\begin{chapterexercise}[Mass and weight]\difficulty{1}
Calculate the weight of a $70\,\mathrm{kg}$ person on Earth using
$g=9.81\,\mathrm{m\,s^{-2}}$.
\end{chapterexercise}

\subsection*{Intermediate}
\begin{chapterexercise}[Inclined plane]\difficulty{2}
A block slides without friction down an incline of angle $\theta$.
Draw the free-body diagram and derive its acceleration along the plane.
\end{chapterexercise}

\subsection*{Advanced}
\begin{chapterexercise}[Momentum form]\difficulty{3}
Starting from $\mathbf F=d\mathbf p/dt$, explain why
$\mathbf F=m\mathbf a$ requires constant mass.
\end{chapterexercise}

\subsection*{Challenge}
\begin{chapterexercise}[Magnetic force preview]\difficulty{4}
The magnetic force is $\mathbf F=q\mathbf v\times\mathbf B$.
Explain why it can change the direction of velocity without changing the
particle's speed.
\end{chapterexercise}

\chapterlookaheadheading
\begin{lookingaheadbox}
The next stage identifies the forces that enter Newton's equation. Of special
importance is the Lorentz force,
\[
\mathbf F=q(\mathbf E+\mathbf v\times\mathbf B),
\]
which later leads to cyclotron motion, Landau quantization, and the Quantum
Hall Effect.
\end{lookingaheadbox}



%%%%%%%%%%%%%%%%%%%%%%%%%%%%%%%%%%%%%%%%%%%%%%%%%%%%%%%%%%%%%%%%%%%%%%%%%%%%%%%%
% Visual language preview for later volumes
%%%%%%%%%%%%%%%%%%%%%%%%%%%%%%%%%%%%%%%%%%%%%%%%%%%%%%%%%%%%%%%%%%%%%%%%%%%%%%%%
\clearpage
\chapter*{Visual Language Preview for Quantum Mechanics}
\addcontentsline{toc}{chapter}{Visual Language Preview}

This unnumbered design page fixes the visual conventions that later chapters
will use. It is a template, not a substitute for the full physics discussion.

\section*{Wavefunctions}

\begin{bookfigure}{Template for one-dimensional wavefunctions. Solid and dashed
curves may represent different stationary states, components, or comparison
solutions, provided each use is explicitly identified in the accompanying text.}
\wavefunctionplot
\end{bookfigure}

\section*{Quantum Hall Geometry}

\begin{bookfigure}{Template for a Hall bar or two-dimensional electron gas.
The magnetic field is perpendicular to the plane; edge-state arrows indicate
opposite propagation directions on opposite boundaries.}
\quantumhalltemplate
\end{bookfigure}

\begin{visualguidebox}[Future consistency rules]
Later quantum figures should preserve a fixed coordinate convention:
the sample lies in the $xy$ plane, the magnetic field points along $z$,
source--drain transport is labelled explicitly, and bulk states are visually
distinguished from boundary-localized edge states.
\end{visualguidebox}


\end{document}

